\documentclass[8pt]{beamer}
\usetheme{Madrid}
\usecolortheme{default}
\usepackage{amsmath, amssymb, graphicx, array, multirow, booktabs, xcolor, tikz}
\usetikzlibrary{shapes,arrows,positioning,fit,backgrounds}

% Compact font settings
\setbeamerfont{title}{size=\Large}
\setbeamerfont{subtitle}{size=\small}
\setbeamerfont{author}{size=\scriptsize}
\setbeamerfont{date}{size=\scriptsize}
\setbeamerfont{frametitle}{size=\large}
\setbeamerfont{framesubtitle}{size=\normalsize}
\setbeamerfont{enumerate item}{size=\footnotesize}
\setbeamerfont{itemize item}{size=\footnotesize}
\setbeamerfont{block title}{size=\footnotesize}

% Compact spacing
\setlength{\itemsep}{0.08ex}
\setlength{\parskip}{0.08ex}
\setlength{\leftmargini}{0.3cm}
\setbeamersize{text margin left=3mm, text margin right=3mm}
\addtobeamertemplate{frame begin}{\vspace{-0.4cm}}{}

% Colors
\definecolor{myblue}{RGB}{0,0,255}
\definecolor{myred}{RGB}{255,0,0}
\definecolor{mygreen}{RGB}{0,128,0}
\definecolor{myorange}{RGB}{255,165,0}
\definecolor{mypurple}{RGB}{128,0,128}
\definecolor{mygray}{RGB}{128,128,128}
\definecolor{mygold}{RGB}{255,215,0}

\title{Supervised Learning for Numerical Data}
\subtitle{100 Questions: Linear Regression, Logistic Regression, LDA}
\author{GARP RAI Exam Preparation}
\date{}

\begin{document}
	
	\begin{frame}
		\titlepage
	\end{frame}
	
	% ============================================================
	% SECTION 1: INTRODUCTION TO LINEAR REGRESSION - QUESTIONS 1-10
	% ============================================================
	
	\begin{frame}{Q1: Simple Linear Regression Definition - Tutorial}
		\fontsize{6.5}{7.5}\selectfont
		\textbf{QUESTION: What is simple linear regression?}
		
		\vspace{0.1cm}
		\textbf{OPTIONS:}
		\begin{enumerate}
			\item A model with multiple features and one target
			\item A model with one feature and one continuous target variable
			\item A model with no features
			\item A model for classification problems
		\end{enumerate}
		
		\vspace{0.1cm}
		\textcolor{myblue}{\textbf{ANSWER: B}}
		
		\vspace{0.1cm}
		\textbf{DETAILED EXPLANATION:}
		\begin{itemize}
			\item Simple linear regression is also called bivariate regression
			\item It involves just two variables: one target (y) and one feature (x)
			\item Models how y changes in response to changes in x
			\item Example: Stock return (y) influenced by market return (x)
		\end{itemize}
		
		\vspace{0.1cm}
		\textcolor{myred}{\textbf{WHY OTHER OPTIONS ARE WRONG:}}
		\begin{itemize}
			\item \textcolor{myred}{A:} This describes MULTIPLE linear regression
			\item \textcolor{myred}{C:} A model with no features would just be the mean
			\item \textcolor{myred}{D:} Linear regression is for continuous targets, not classification
		\end{itemize}
		
		\textcolor{myred}{\textbf{EXAM TRAP:}} Simple linear regression = \textcolor{myblue}{one feature + one continuous target}!
	\end{frame}
	


	\begin{frame}{Q2: Linear Regression Equation - Tutorial}
	\fontsize{6.5}{7.5}\selectfont
	\textbf{QUESTION: What is the equation for a simple linear regression model?}
	
	\vspace{0.1cm}
	\textbf{OPTIONS:}
	\begin{enumerate}
		\item $y_i = \beta_0 + \beta_1 x_i + u_i$
		\item $y_i = \beta_0 + \beta_1 x_{1i} + \beta_2 x_{2i} + u_i$
		\item $y_i = \beta_0 + u_i$
		\item $y_i = \beta_1 x_i + u_i$
	\end{enumerate}
	
	\vspace{0.1cm}
	\textcolor{myblue}{\textbf{ANSWER: A}}
	
	\vspace{0.1cm}
	\textbf{DETAILED EXPLANATION:}
	\begin{itemize}
		\item $y_i$: Target variable (dependent variable)
		\item $x_i$: Feature (independent variable)
		\item $\beta_0$: Intercept (bias) - value of y when x=0
		\item $\beta_1$: Slope (weight) - change in y for 1-unit change in x
		\item $u_i$: Error term - captures unexplained variation
	\end{itemize}
	
	\vspace{0.1cm}
	\textcolor{myred}{\textbf{WHY OTHER OPTIONS ARE WRONG:}}
	\begin{itemize}
		\item \textcolor{myred}{B:} This is MULTIPLE linear regression
		\item \textcolor{myred}{C:} Missing the feature term
		\item \textcolor{myred}{D:} Missing the intercept term
	\end{itemize}
	
	\textcolor{myred}{\textbf{EXAM TRAP:}} Simple linear regression has \textcolor{myblue}{intercept + slope + error}!
\end{frame}


	
	\begin{frame}{Q3: Error Term Definition - Tutorial}
		\fontsize{6.5}{7.5}\selectfont
		\textbf{QUESTION: What does the error term $u_i$ represent in a regression model?}
		
		\vspace{0.1cm}
		\textbf{OPTIONS:}
		\begin{enumerate}
			\item The predicted value
			\item The intercept of the model
			\item Unexplained factors affecting y not in the model
			\item The slope of the model
		\end{enumerate}
		
		\vspace{0.1cm}
		\textcolor{myblue}{\textbf{ANSWER: C}}
		
		\vspace{0.1cm}
		\textbf{DETAILED EXPLANATION:}
		\begin{itemize}
			\item Captures all other factors affecting y
			\item Includes inherent randomness
			\item Assumed to have mean of zero
			\item Assumed to have constant variance
			\item Assumed to be independent
		\end{itemize}
		
		\vspace{0.1cm}
		\textcolor{myred}{\textbf{WHY OTHER OPTIONS ARE WRONG:}}
		\begin{itemize}
			\item \textcolor{myred}{A:} Predicted value is $\hat{y}_i$
			\item \textcolor{myred}{B:} Intercept is $\beta_0$
			\item \textcolor{myred}{D:} Slope is $\beta_1$
		\end{itemize}
		
		\textcolor{myred}{\textbf{EXAM TRAP:}} Error term captures \textcolor{myblue}{unexplained variation}!
	\end{frame}
	
	\begin{frame}{Q4: Parameter Estimation Methods - Tutorial}
		\fontsize{6.5}{7.5}\selectfont
		\textbf{QUESTION: What are the three primary methods for estimating regression parameters?}
		
		\vspace{0.1cm}
		\textbf{OPTIONS:}
		\begin{enumerate}
			\item Least Squares, Maximum Likelihood, Method of Moments
			\item Gradient Descent, Backpropagation, Cross-Validation
			\item LASSO, Ridge, Elastic Net
			\item Linear, Logistic, Polynomial
		\end{enumerate}
		
		\vspace{0.1cm}
		\textcolor{myblue}{\textbf{ANSWER: A}}
		
		\vspace{0.1cm}
		\textbf{DETAILED EXPLANATION:}
		\begin{itemize}
			\item \textbf{Least Squares:} Most common (OLS)
			\item \textbf{Maximum Likelihood:} More general method
			\item \textbf{Method of Moments:} Another statistical technique
			\item OLS is most widely used for linear regression
		\end{itemize}
		
		\vspace{0.1cm}
		\textcolor{myred}{\textbf{WHY OTHER OPTIONS ARE WRONG:}}
		\begin{itemize}
			\item \textcolor{myred}{B:} These are optimization/ML techniques
			\item \textcolor{myred}{C:} These are regularization methods
			\item \textcolor{myred}{D:} These are model types, not estimation methods
		\end{itemize}
		
		\textcolor{myred}{\textbf{EXAM TRAP:}} Three estimation methods = \textcolor{myblue}{OLS, MLE, Method of Moments}!
	\end{frame}
	
	\begin{frame}{Q5: OLS Definition - Tutorial}
		\fontsize{6.5}{7.5}\selectfont
		\textbf{QUESTION: What is Ordinary Least Squares (OLS)?}
		
		\vspace{0.1cm}
		\textbf{OPTIONS:}
		\begin{enumerate}
			\item A method that maximizes the likelihood of the data
			\item A method that minimizes the sum of squared differences between actual and predicted values
			\item A method that minimizes absolute errors
			\item A method that randomly selects parameters
		\end{enumerate}
		
		\vspace{0.1cm}
		\textcolor{myblue}{\textbf{ANSWER: B}}
		
		\vspace{0.1cm}
		\textbf{DETAILED EXPLANATION:}
		\begin{itemize}
			\item Most common estimation method for linear regression
			\item Finds $\beta_0$ and $\beta_1$ that minimize SSE
			\item SSE = $\sum (y_i - \hat{y}_i)^2$
			\item Squaring penalizes large errors more heavily
		\end{itemize}
		
		\vspace{0.1cm}
		\textcolor{myred}{\textbf{WHY OTHER OPTIONS ARE WRONG:}}
		\begin{itemize}
			\item \textcolor{myred}{A:} This describes Maximum Likelihood Estimation
			\item \textcolor{myred}{C:} This describes Least Absolute Deviations (LAD)
			\item \textcolor{myred}{D:} This describes random search
		\end{itemize}
		
		\textcolor{myred}{\textbf{EXAM TRAP:}} OLS minimizes \textcolor{myblue}{sum of squared errors}!
	\end{frame}
	
	\begin{frame}{Q6: Multiple Linear Regression - Tutorial}
		\fontsize{6.5}{7.5}\selectfont
		\textbf{QUESTION: What is multiple linear regression?}
		
		\vspace{0.1cm}
		\textbf{OPTIONS:}
		\begin{enumerate}
			\item A model with one feature and one target
			\item A model with multiple features and a continuous target
			\item A model for classification problems
			\item A model with no features
		\end{enumerate}
		
		\vspace{0.1cm}
		\textcolor{myblue}{\textbf{ANSWER: B}}
		
		\vspace{0.1cm}
		\textbf{DETAILED EXPLANATION:}
		\begin{itemize}
			\item Direct extension of simple linear regression
			\item Includes multiple features ($x_1, x_2, ..., x_m$)
			\item Equation: $y_i = \beta_0 + \beta_1 x_{1i} + ... + \beta_m x_{mi} + u_i$
			\item Each coefficient measures partial effect
		\end{itemize}
		
		\vspace{0.1cm}
		\textcolor{myred}{\textbf{WHY OTHER OPTIONS ARE WRONG:}}
		\begin{itemize}
			\item \textcolor{myred}{A:} This describes SIMPLE linear regression
			\item \textcolor{myred}{C:} Multiple linear regression is for continuous targets
			\item \textcolor{myred}{D:} A model with no features is just the mean
		\end{itemize}
		
		\textcolor{myred}{\textbf{EXAM TRAP:}} Multiple linear regression = \textcolor{myblue}{multiple features + continuous target}!
	\end{frame}
	
	\begin{frame}{Q7: Multiple Linear Regression Equation - Tutorial}
		\fontsize{6.5}{7.5}\selectfont
		\textbf{QUESTION: What is the equation for multiple linear regression with m features?}
		
		\vspace{0.1cm}
		\textbf{OPTIONS:}
		\begin{enumerate}
			\item $y_i = \beta_0 + \beta_1 x_i + u_i$
			\item $y_i = \beta_0 + \beta_1 x_{1i} + \beta_2 x_{2i} + ... + \beta_m x_{mi} + u_i$
			\item $y_i = \beta_0 + u_i$
			\item $y_i = \beta_1 x_{1i} + \beta_2 x_{2i} + u_i$
		\end{enumerate}
		
		\vspace{0.1cm}
		\textcolor{myblue}{\textbf{ANSWER: B}}
		
		\vspace{0.1cm}
		\textbf{DETAILED EXPLANATION:}
		\begin{itemize}
			\item $m$ features require $m+1$ parameters
			\item $\beta_0$: intercept
			\item $\beta_1, ..., \beta_m$: slope parameters
			\item Each $\beta$ measures partial effect
			\item OLS can still be used for estimation
		\end{itemize}
		
		\vspace{0.1cm}
		\textcolor{myred}{\textbf{WHY OTHER OPTIONS ARE WRONG:}}
		\begin{itemize}
			\item \textcolor{myred}{A:} This is SIMPLE linear regression
			\item \textcolor{myred}{C:} Missing all feature terms
			\item \textcolor{myred}{D:} Missing the intercept
		\end{itemize}
		
		\textcolor{myred}{\textbf{EXAM TRAP:}} Multiple regression has \textcolor{myblue}{m features + intercept} = m+1 parameters!
	\end{frame}
	
	\begin{frame}{Q8: Partial Effect Interpretation - Tutorial}
		\fontsize{6.5}{7.5}\selectfont
		\textbf{QUESTION: In multiple regression, what does each $\beta$ coefficient measure?}
		
		\vspace{0.1cm}
		\textbf{OPTIONS:}
		\begin{enumerate}
			\item The total effect of a feature
			\item The partial effect of a feature controlling for other features
			\item The correlation between features
			\item The error in the model
		\end{enumerate}
		
		\vspace{0.1cm}
		\textcolor{myblue}{\textbf{ANSWER: B}}
		
		\vspace{0.1cm}
		\textbf{DETAILED EXPLANATION:}
		\begin{itemize}
			\item Measures effect of one feature while holding others constant
			\item "Ceteris paribus" interpretation
			\item Controls for other variables in the model
			\item Example: Effect of experience on salary holding degree constant
		\end{itemize}
		
		\vspace{0.1cm}
		\textcolor{myred}{\textbf{WHY OTHER OPTIONS ARE WRONG:}}
		\begin{itemize}
			\item \textcolor{myred}{A:} Total effect is different from partial effect
			\item \textcolor{myred}{C:} Correlation is a different concept
			\item \textcolor{myred}{D:} Error is $u_i$, not $\beta$
		\end{itemize}
		
		\textcolor{myred}{\textbf{EXAM TRAP:}} $\beta$ measures \textcolor{myblue}{partial effect} controlling for other variables!
	\end{frame}
	
	\begin{frame}{Q9: Interaction Term - Tutorial}
		\fontsize{6.5}{7.5}\selectfont
		\textbf{QUESTION: What is an interaction term in a regression model?}
		
		\vspace{0.1cm}
		\textbf{OPTIONS:}
		\begin{enumerate}
			\item $x_1 + x_2$ where features are added
			\item $x_1 \times x_2$ where features are multiplied
			\item $x_1^2$ where a feature is squared
			\item $\log(x_1)$ where a feature is transformed
		\end{enumerate}
		
		\vspace{0.1cm}
		\textcolor{myblue}{\textbf{ANSWER: B}}
		
		\vspace{0.1cm}
		\textbf{DETAILED EXPLANATION:}
		\begin{itemize}
			\item Equation: $y_i = \beta_0 + \beta_1 x_{1i} + \beta_2 x_{2i} + \beta_3 x_{1i} x_{2i} + u_i$
			\item Captures complementarity between features
			\item Effect of $x_1$ depends on $x_2$ and vice versa
			\item Changes in both variables together impact y
		\end{itemize}
		
		\vspace{0.1cm}
		\textcolor{myred}{\textbf{WHY OTHER OPTIONS ARE WRONG:}}
		\begin{itemize}
			\item \textcolor{myred}{A:} This is addition, not interaction
			\item \textcolor{myred}{C:} This is a quadratic term
			\item \textcolor{myred}{D:} This is a log transformation
		\end{itemize}
		
		\textcolor{myred}{\textbf{EXAM TRAP:}} Interaction term = \textcolor{myblue}{features multiplied together}!
	\end{frame}
	
	\begin{frame}{Q10: Quadratic Regression - Tutorial}
		\fontsize{6.5}{7.5}\selectfont
		\textbf{QUESTION: What does adding a squared term to a regression allow?}
		
		\vspace{0.1cm}
		\textbf{OPTIONS:}
		\begin{enumerate}
			\item A linear relationship
			\item A U-shaped or inverted U-shaped relationship
			\item No relationship between variables
			\item A perfectly flat relationship
		\end{enumerate}
		
		\vspace{0.1cm}
		\textcolor{myblue}{\textbf{ANSWER: B}}
		
		\vspace{0.1cm}
		\textbf{DETAILED EXPLANATION:}
		\begin{itemize}
			\item Equation: $y_i = \beta_0 + \beta_1 x_i + \beta_2 x_i^2 + u_i$
			\item Allows for nonlinear relationships
			\item U-shaped ($\beta_2 > 0$): decreases then increases
			\item Inverted U-shaped ($\beta_2 < 0$): increases then decreases
			\item Example: Experience and salary (diminishing returns)
		\end{itemize}
		
		\vspace{0.1cm}
		\textcolor{myred}{\textbf{WHY OTHER OPTIONS ARE WRONG:}}
		\begin{itemize}
			\item \textcolor{myred}{A:} Linear relationship is without squared term
			\item \textcolor{myred}{C:} There is still a relationship
			\item \textcolor{myred}{D:} Flat relationship would require $\beta_1 = \beta_2 = 0$
		\end{itemize}
		
		\textcolor{myred}{\textbf{EXAM TRAP:}} Quadratic term allows \textcolor{myblue}{U-shaped relationships}!
	\end{frame}
	
	% ============================================================
	% SECTION 2: REGRESSION PROBLEMS - QUESTIONS 11-25
	% ============================================================
	
	\begin{frame}{Q11: Omitting Relevant Features - Tutorial}
		\fontsize{6.5}{7.5}\selectfont
		\textbf{QUESTION: What is the consequence of omitting a relevant feature from a regression?}
		
		\vspace{0.1cm}
		\textbf{OPTIONS:}
		\begin{enumerate}
			\item No consequence
			\item Biased parameter estimates for included features
			\item Smaller standard errors
			\item Better model fit
		\end{enumerate}
		
		\vspace{0.1cm}
		\textcolor{myblue}{\textbf{ANSWER: B}}
		
		\vspace{0.1cm}
		\textbf{DETAILED EXPLANATION:}
		\begin{itemize}
			\item Omitted variable bias
			\item Estimates are systematically wrong
			\item Bias does not disappear with larger samples
			\item Model attributes influence of missing variable to included ones
		\end{itemize}
		
		\vspace{0.1cm}
		\textcolor{myred}{\textbf{WHY OTHER OPTIONS ARE WRONG:}}
		\begin{itemize}
			\item \textcolor{myred}{A:} There are significant consequences
			\item \textcolor{myred}{C:} Standard errors may be biased too
			\item \textcolor{myred}{D:} Model fit is compromised
		\end{itemize}
		
		\textcolor{myred}{\textbf{EXAM TRAP:}} Omitting relevant features causes \textcolor{myblue}{biased estimates}!
	\end{frame}
	
	\begin{frame}{Q12: Including Irrelevant Features - Tutorial}
		\fontsize{6.5}{7.5}\selectfont
		\textbf{QUESTION: What is the consequence of including an irrelevant feature?}
		
		\vspace{0.1cm}
		\textbf{OPTIONS:}
		\begin{enumerate}
			\item Biased estimates for relevant features
			\item Inefficient (less precise) parameter estimates
			\item No consequence at all
			\item Smaller standard errors
		\end{enumerate}
		
		\vspace{0.1cm}
		\textcolor{myblue}{\textbf{ANSWER: B}}
		
		\vspace{0.1cm}
		\textbf{DETAILED EXPLANATION:}
		\begin{itemize}
			\item Typically not biased for relevant variables
			\item Causes inefficiency - larger standard errors
			\item Model becomes overly complex
			\item Risk of overfitting
		\end{itemize}
		
		\vspace{0.1cm}
		\textcolor{myred}{\textbf{WHY OTHER OPTIONS ARE WRONG:}}
		\begin{itemize}
			\item \textcolor{myred}{A:} Relevant estimates are usually unbiased
			\item \textcolor{myred}{C:} There is a consequence (inefficiency)
			\item \textcolor{myred}{D:} Standard errors increase, not decrease
		\end{itemize}
		
		\textcolor{myred}{\textbf{EXAM TRAP:}} Irrelevant features cause \textcolor{myblue}{inefficiency} (larger standard errors)!
	\end{frame}
	
	\begin{frame}{Q13: Incorrect Functional Form - Tutorial}
		\fontsize{6.5}{7.5}\selectfont
		\textbf{QUESTION: What happens if we assume a linear relationship when the true relationship is non-linear?}
		
		\vspace{0.1cm}
		\textbf{OPTIONS:}
		\begin{enumerate}
			\item The model will fit perfectly
			\item The model will be a poor representation of reality
			\item The model will have no errors
			\item The model will be more efficient
		\end{enumerate}
		
		\vspace{0.1cm}
		\textcolor{myblue}{\textbf{ANSWER: B}}
		
		\vspace{0.1cm}
		\textbf{DETAILED EXPLANATION:}
		\begin{itemize}
			\item Model misspecification
			\item Biased predictions
			\item Incorrect inferences about feature impacts
			\item Example: Linear model for quadratic relationship
		\end{itemize}
		
		\vspace{0.1cm}
		\textcolor{myred}{\textbf{WHY OTHER OPTIONS ARE WRONG:}}
		\begin{itemize}
			\item \textcolor{myred}{A:} Misspecified models don't fit perfectly
			\item \textcolor{myred}{C:} There will be systematic errors
			\item \textcolor{myred}{D:} Efficiency is compromised
		\end{itemize}
		
		\textcolor{myred}{\textbf{EXAM TRAP:}} Wrong functional form = \textcolor{myblue}{poor model representation}!
	\end{frame}
	
	\begin{frame}{Q14: Multicollinearity Definition - Tutorial}
		\fontsize{6.5}{7.5}\selectfont
		\textbf{QUESTION: What is multicollinearity in regression?}
		
		\vspace{0.1cm}
		\textbf{OPTIONS:}
		\begin{enumerate}
			\item When features are highly correlated with the target
			\item When features are highly correlated with each other
			\item When the error term has non-constant variance
			\item When there are outliers in the data
		\end{enumerate}
		
		\vspace{0.1cm}
		\textcolor{myblue}{\textbf{ANSWER: B}}
		
		\vspace{0.1cm}
		\textbf{DETAILED EXPLANATION:}
		\begin{itemize}
			\item Features (independent variables) are highly correlated
			\item Makes it hard to isolate individual effects
			\item Two types: perfect and imperfect
			\item Example: Square footage and number of bedrooms
		\end{itemize}
		
		\vspace{0.1cm}
		\textcolor{myred}{\textbf{WHY OTHER OPTIONS ARE WRONG:}}
		\begin{itemize}
			\item \textcolor{myred}{A:} Correlation with target is desirable
			\item \textcolor{myred}{C:} This is heteroskedasticity
			\item \textcolor{myred}{D:} This is a different problem
		\end{itemize}
		
		\textcolor{myred}{\textbf{EXAM TRAP:}} Multicollinearity = \textcolor{myblue}{features correlated with each other}!
	\end{frame}
	
	\begin{frame}{Q15: Perfect Multicollinearity - Tutorial}
		\fontsize{6.5}{7.5}\selectfont
		\textbf{QUESTION: What happens with perfect multicollinearity?}
		
		\vspace{0.1cm}
		\textbf{OPTIONS:}
		\begin{enumerate}
			\item OLS estimation works normally
			\item OLS estimation breaks down; separate effects cannot be estimated
			\item Standard errors become smaller
			\item The model becomes more efficient
		\end{enumerate}
		
		\vspace{0.1cm}
		\textcolor{myblue}{\textbf{ANSWER: B}}
		
		\vspace{0.1cm}
		\textbf{DETAILED EXPLANATION:}
		\begin{itemize}
			\item Exact linear relationship between features
			\item Example: Celsius and Fahrenheit included together
			\item OLS cannot estimate separate coefficients
			\item Must remove one or more features
		\end{itemize}
		
		\vspace{0.1cm}
		\textcolor{myred}{\textbf{WHY OTHER OPTIONS ARE WRONG:}}
		\begin{itemize}
			\item \textcolor{myred}{A:} OLS fails with perfect multicollinearity
			\item \textcolor{myred}{C:} Standard errors are not valid
			\item \textcolor{myred}{D:} Model is not efficient
		\end{itemize}
		
		\textcolor{myred}{\textbf{EXAM TRAP:}} Perfect multicollinearity = \textcolor{myblue}{OLS breaks down}!
	\end{frame}
	
	\begin{frame}{Q16: Imperfect Multicollinearity - Tutorial}
		\fontsize{6.5}{7.5}\selectfont
		\textbf{QUESTION: What is imperfect multicollinearity?}
		
		\vspace{0.1cm}
		\textbf{OPTIONS:}
		\begin{enumerate}
			\item Features are exactly linearly related
			\item Features are very closely (but not perfectly) correlated
			\item Features are completely uncorrelated
			\item Features are correlated with the error term
		\end{enumerate}
		
		\vspace{0.1cm}
		\textcolor{myblue}{\textbf{ANSWER: B}}
		
		\vspace{0.1cm}
		\textbf{DETAILED EXPLANATION:}
		\begin{itemize}
			\item Correlation very high (e.g., 0.9 or -0.95)
			\item More common than perfect collinearity
			\item Causes large standard errors
			\item Makes it hard to separate effects
		\end{itemize}
		
		\vspace{0.1cm}
		\textcolor{myred}{\textbf{WHY OTHER OPTIONS ARE WRONG:}}
		\begin{itemize}
			\item \textcolor{myred}{A:} This is perfect multicollinearity
			\item \textcolor{myred}{C:} This is no multicollinearity
			\item \textcolor{myred}{D:} This is endogeneity, not multicollinearity
		\end{itemize}
		
		\textcolor{myred}{\textbf{EXAM TRAP:}} Imperfect multicollinearity = \textcolor{myblue}{very high but not perfect correlation}!
	\end{frame}
	
	\begin{frame}{Q17: Outlier Definition - Tutorial}
		\fontsize{6.5}{7.5}\selectfont
		\textbf{QUESTION: What is an outlier in regression?}
		
		\vspace{0.1cm}
		\textbf{OPTIONS:}
		\begin{enumerate}
			\item A data point that fits the pattern perfectly
			\item An anomalous data point far from the bulk of observations
			\item The mean of the data
			\item The median of the data
		\end{enumerate}
		
		\vspace{0.1cm}
		\textcolor{myblue}{\textbf{ANSWER: B}}
		
		\vspace{0.1cm}
		\textbf{DETAILED EXPLANATION:}
		\begin{itemize}
			\item Data point with large residual
			\item Far from the regression line
			\item Can be due to data entry errors
			\item Or rare but genuine events
		\end{itemize}
		
		\vspace{0.1cm}
		\textcolor{myred}{\textbf{WHY OTHER OPTIONS ARE WRONG:}}
		\begin{itemize}
			\item \textcolor{myred}{A:} This is an inlier, not an outlier
			\item \textcolor{myred}{C:} The mean is a summary statistic
			\item \textcolor{myred}{D:} The median is a summary statistic
		\end{itemize}
		
		\textcolor{myred}{\textbf{EXAM TRAP:}} Outliers are \textcolor{myblue}{anomalous data points far from others}!
	\end{frame}
	
	\begin{frame}{Q18: Outlier Impact on OLS - Tutorial}
		\fontsize{6.5}{7.5}\selectfont
		\textbf{QUESTION: Why can outliers have a large impact on OLS estimates?}
		
		\vspace{0.1cm}
		\textbf{OPTIONS:}
		\begin{enumerate}
			\item OLS ignores outliers
			\item OLS minimizes squared residuals, so large residuals have disproportionate influence
			\item OLS gives equal weight to all observations
			\item OLS uses absolute residuals
		\end{enumerate}
		
		\vspace{0.1cm}
		\textcolor{myblue}{\textbf{ANSWER: B}}
		
		\vspace{0.1cm}
		\textbf{DETAILED EXPLANATION:}
		\begin{itemize}
			\item Squaring residuals amplifies large differences
			\item Points far from line pull the line toward them
			\item Can significantly change parameter estimates
			\item Example: One machine with many failures
		\end{itemize}
		
		\vspace{0.1cm}
		\textcolor{myred}{\textbf{WHY OTHER OPTIONS ARE WRONG:}}
		\begin{itemize}
			\item \textcolor{myred}{A:} OLS does NOT ignore outliers
			\item \textcolor{myred}{C:} Squaring gives more weight to large residuals
			\item \textcolor{myred}{D:} OLS uses squared, not absolute residuals
		\end{itemize}
		
		\textcolor{myred}{\textbf{EXAM TRAP:}} Outliers have \textcolor{myblue}{disproportionate influence} because of squared errors!
	\end{frame}
	
	\begin{frame}{Q19: Cook's Distance - Tutorial}
		\fontsize{6.5}{7.5}\selectfont
		\textbf{QUESTION: What does Cook's Distance measure?}
		
		\vspace{0.1cm}
		\textbf{OPTIONS:}
		\begin{enumerate}
			\item The mean of the residuals
			\item The effect of deleting an observation on all fitted values
			\item The correlation between features
			\item The variance of the error term
		\end{enumerate}
		
		\vspace{0.1cm}
		\textcolor{myblue}{\textbf{ANSWER: B}}
		
		\vspace{0.1cm}
		\textbf{DETAILED EXPLANATION:}
		\begin{itemize}
			\item Diagnostic for influential points
			\item Measures how much fitted values change when point is removed
			\item Large values indicate influential observations
			\item Higher than others warrants investigation
		\end{itemize}
		
		\vspace{0.1cm}
		\textcolor{myred}{\textbf{WHY OTHER OPTIONS ARE WRONG:}}
		\begin{itemize}
			\item \textcolor{myred}{A:} Mean of residuals is different
			\item \textcolor{myred}{C:} This is correlation
			\item \textcolor{myred}{D:} This is error variance
		\end{itemize}
		
		\textcolor{myred}{\textbf{EXAM TRAP:}} Cook's Distance identifies \textcolor{myblue}{influential observations}!
	\end{frame}
	
	\begin{frame}{Q20: Heteroskedasticity Definition - Tutorial}
		\fontsize{6.5}{7.5}\selectfont
		\textbf{QUESTION: What is heteroskedasticity in regression?}
		
		\vspace{0.1cm}
		\textbf{OPTIONS:}
		\begin{enumerate}
			\item Constant variance of errors
			\item Non-constant variance of errors
			\item Zero mean of errors
			\item Independence of errors
		\end{enumerate}
		
		\vspace{0.1cm}
		\textcolor{myblue}{\textbf{ANSWER: B}}
		
		\vspace{0.1cm}
		\textbf{DETAILED EXPLANATION:}
		\begin{itemize}
			\item Error variance is not constant
			\item Opposite of homoskedasticity
			\item Common in cross-sectional data
			\item Example: Profits of small vs large companies
		\end{itemize}
		
		\vspace{0.1cm}
		\textcolor{myred}{\textbf{WHY OTHER OPTIONS ARE WRONG:}}
		\begin{itemize}
			\item \textcolor{myred}{A:} This is homoskedasticity
			\item \textcolor{myred}{C:} This is an OLS assumption
			\item \textcolor{myred}{D:} This is another OLS assumption
		\end{itemize}
		
		\textcolor{myred}{\textbf{EXAM TRAP:}} Heteroskedasticity = \textcolor{myblue}{non-constant error variance}!
	\end{frame}
	
	\begin{frame}{Q21: Heteroskedasticity Detection - Tutorial}
		\fontsize{6.5}{7.5}\selectfont
		\textbf{QUESTION: How can heteroskedasticity be detected?}
		
		\vspace{0.1cm}
		\textbf{OPTIONS:}
		\begin{enumerate}
			\item Looking at residual plots for changing spread
			\item Looking at the mean of residuals
			\item Looking at the R-squared
			\item Looking at the intercept
		\end{enumerate}
		
		\vspace{0.1cm}
		\textcolor{myblue}{\textbf{ANSWER: A}}
		
		\vspace{0.1cm}
		\textbf{DETAILED EXPLANATION:}
		\begin{itemize}
			\item Plot residuals vs fitted values
			\item Look for funnel shape
			\item Spread changes with fitted values
			\item Formal test: White's Test
		\end{itemize}
		
		\vspace{0.1cm}
		\textcolor{myred}{\textbf{WHY OTHER OPTIONS ARE WRONG:}}
		\begin{itemize}
			\item \textcolor{myred}{B:} Mean of residuals is usually zero
			\item \textcolor{myred}{C:} R-squared doesn't detect heteroskedasticity
			\item \textcolor{myred}{D:} Intercept doesn't detect heteroskedasticity
		\end{itemize}
		
		\textcolor{myred}{\textbf{EXAM TRAP:}} Heteroskedasticity = \textcolor{myblue}{funnel shape in residual plot}!
	\end{frame}
	
	\begin{frame}{Q22: White's Test - Tutorial}
		\fontsize{6.5}{7.5}\selectfont
		\textbf{QUESTION: What is White's Test used for?}
		
		\vspace{0.1cm}
		\textbf{OPTIONS:}
		\begin{enumerate}
			\item Detecting multicollinearity
			\item Detecting heteroskedasticity
			\item Detecting outliers
			\item Detecting autocorrelation
		\end{enumerate}
		
		\vspace{0.1cm}
		\textcolor{myblue}{\textbf{ANSWER: B}}
		
		\vspace{0.1cm}
		\textbf{DETAILED EXPLANATION:}
		\begin{itemize}
			\item Formal statistical test for heteroskedasticity
			\item Standard option in most regression software
			\item Significant result indicates heteroskedasticity
			\item Produces White's heteroskedasticity-consistent errors
		\end{itemize}
		
		\vspace{0.1cm}
		\textcolor{myred}{\textbf{WHY OTHER OPTIONS ARE WRONG:}}
		\begin{itemize}
			\item \textcolor{myred}{A:} Variance Inflation Factor (VIF) detects multicollinearity
			\item \textcolor{myred}{C:} Cook's Distance detects outliers
			\item \textcolor{myred}{D:} Durbin-Watson detects autocorrelation
		\end{itemize}
		
		\textcolor{myred}{\textbf{EXAM TRAP:}} White's Test = \textcolor{myblue}{heteroskedasticity detection}!
	\end{frame}
	
	\begin{frame}{Q23: Heteroskedasticity Remedies - Tutorial}
		\fontsize{6.5}{7.5}\selectfont
		\textbf{QUESTION: What is one remedy for heteroskedasticity?}
		
		\vspace{0.1cm}
		\textbf{OPTIONS:}
		\begin{enumerate}
			\item Add more features
			\item Use Weighted Least Squares or robust standard errors
			\item Remove all features
			\item Ignore it
		\end{enumerate}
		
		\vspace{0.1cm}
		\textcolor{myblue}{\textbf{ANSWER: B}}
		
		\vspace{0.1cm}
		\textbf{DETAILED EXPLANATION:}
		\begin{itemize}
			\item \textbf{WLS:} Weight observations by inverse variance
			\item \textbf{Robust errors:} White's heteroskedasticity-consistent errors
			\item \textbf{Log transformation:} Of variables can help
			\item Select "robust errors" option in software
		\end{itemize}
		
		\vspace{0.1cm}
		\textcolor{myred}{\textbf{WHY OTHER OPTIONS ARE WRONG:}}
		\begin{itemize}
			\item \textcolor{myred}{A:} Adding features doesn't fix heteroskedasticity
			\item \textcolor{myred}{C:} Removing all features removes the model
			\item \textcolor{myred}{D:} Ignoring it is not a solution
		\end{itemize}
		
		\textcolor{myred}{\textbf{EXAM TRAP:}} Remedies = \textcolor{myblue}{WLS, robust errors, or log transformation}!
	\end{frame}
	
	\begin{frame}{Q24: Stepwise Regression Purpose - Tutorial}
		\fontsize{6.5}{7.5}\selectfont
		\textbf{QUESTION: What is the purpose of stepwise regression?}
		
		\vspace{0.1cm}
		\textbf{OPTIONS:}
		\begin{enumerate}
			\item To estimate parameters
			\item To automate feature selection
			\item To create new features
			\item To transform variables
		\end{enumerate}
		
		\vspace{0.1cm}
		\textcolor{myblue}{\textbf{ANSWER: B}}
		
		\vspace{0.1cm}
		\textbf{DETAILED EXPLANATION:}
		\begin{itemize}
			\item Selects subset of features from candidate pool
			\item Reduces model complexity
			\item Avoids redundant features
			\item Uses AIC or other criteria
		\end{itemize}
		
		\vspace{0.1cm}
		\textcolor{myred}{\textbf{WHY OTHER OPTIONS ARE WRONG:}}
		\begin{itemize}
			\item \textcolor{myred}{A:} Parameter estimation is OLS
			\item \textcolor{myred}{C:} Feature creation is different
			\item \textcolor{myred}{D:} Variable transformation is different
		\end{itemize}
		
		\textcolor{myred}{\textbf{EXAM TRAP:}} Stepwise regression = \textcolor{myblue}{automated feature selection}!
	\end{frame}
	
	\begin{frame}{Q25: Forward vs Backward Selection - Tutorial}
		\fontsize{6.5}{7.5}\selectfont
		\textbf{QUESTION: How does forward selection differ from backward selection?}
		
		\vspace{0.1cm}
		\textbf{OPTIONS:}
		\begin{enumerate}
			\item Forward starts with no features; backward starts with all features
			\item Forward starts with all features; backward starts with no features
			\item They are identical
			\item Forward removes features; backward adds features
		\end{enumerate}
		
		\vspace{0.1cm}
		\textcolor{myblue}{\textbf{ANSWER: A}}
		
		\vspace{0.1cm}
		\textbf{DETAILED EXPLANATION:}
		\begin{itemize}
			\item \textbf{Forward:} Start with intercept only, add features one by one
			\item \textbf{Backward:} Start with all features, remove one by one
			\item Both use AIC to decide
			\item May yield different final models
		\end{itemize}
		
		\vspace{0.1cm}
		\textcolor{myred}{\textbf{WHY OTHER OPTIONS ARE WRONG:}}
		\begin{itemize}
			\item \textcolor{myred}{B:} This is reversed
			\item \textcolor{myred}{C:} They are different approaches
			\item \textcolor{myred}{D:} Forward adds; backward removes
		\end{itemize}
		
		\textcolor{myred}{\textbf{EXAM TRAP:}} Forward = \textcolor{myblue}{add features}; Backward = \textcolor{myblue}{remove features}!
	\end{frame}
	
	% ============================================================
	% SECTION 3: LOGISTIC REGRESSION - QUESTIONS 26-45
	% ============================================================
	
	\begin{frame}{Q26: Classification Problem - Tutorial}
		\fontsize{6.5}{7.5}\selectfont
		\textbf{QUESTION: What is a classification problem in machine learning?}
		
		\vspace{0.1cm}
		\textbf{OPTIONS:}
		\begin{enumerate}
			\item Predicting a continuous numerical value
			\item Assigning observations to predefined categories
			\item Clustering similar observations
			\item Reducing dimensionality
		\end{enumerate}
		
		\vspace{0.1cm}
		\textcolor{myblue}{\textbf{ANSWER: B}}
		
		\vspace{0.1cm}
		\textbf{DETAILED EXPLANATION:}
		\begin{itemize}
			\item Target is categorical (not continuous)
			\item Binary: two categories (0/1)
			\item Multi-class: more than two categories
			\item Examples: Default/No Default, Fraud/Legitimate
		\end{itemize}
		
		\vspace{0.1cm}
		\textcolor{myred}{\textbf{WHY OTHER OPTIONS ARE WRONG:}}
		\begin{itemize}
			\item \textcolor{myred}{A:} This is regression
			\item \textcolor{myred}{C:} This is unsupervised learning
			\item \textcolor{myred}{D:} This is feature reduction
		\end{itemize}
		
		\textcolor{myred}{\textbf{EXAM TRAP:}} Classification = \textcolor{myblue}{assigning to categories}!
	\end{frame}
	
	\begin{frame}{Q27: Binary Classification Examples - Tutorial}
		\fontsize{6.5}{7.5}\selectfont
		\textbf{QUESTION: Which is NOT an example of binary classification in risk management?}
		
		\vspace{0.1cm}
		\textbf{OPTIONS:}
		\begin{enumerate}
			\item Loan default (Yes/No)
			\item Fraud detection (Fraudulent/Legitimate)
			\item Predicting house prices
			\item Cybersecurity threat (Malicious/Benign)
		\end{enumerate}
		
		\vspace{0.1cm}
		\textcolor{myblue}{\textbf{ANSWER: C}}
		
		\vspace{0.1cm}
		\textbf{DETAILED EXPLANATION:}
		\begin{itemize}
			\item Binary classification has two categories
			\item Predicting house prices is regression (continuous)
			\item Default prediction, fraud detection, threat detection are all binary
		\end{itemize}
		
		\vspace{0.1cm}
		\textcolor{myred}{\textbf{WHY OTHER OPTIONS ARE WRONG:}}
		\begin{itemize}
			\item \textcolor{myred}{A:} Yes/No default is binary classification
			\item \textcolor{myred}{B:} Fraudulent/Legitimate is binary classification
			\item \textcolor{myred}{D:} Malicious/Benign is binary classification
		\end{itemize}
		
		\textcolor{myred}{\textbf{EXAM TRAP:}} Predicting prices = \textcolor{myblue}{regression}, not classification!
	\end{frame}
	
	\begin{frame}{Q28: Why Linear Regression Fails for Classification - Tutorial}
		\fontsize{6.5}{7.5}\selectfont
		\textbf{QUESTION: Why is linear regression inappropriate for binary classification?}
		
		\vspace{0.1cm}
		\textbf{OPTIONS:}
		\begin{enumerate}
			\item Linear regression is too complex
			\item Predictions can fall outside the 0-1 range
			\item Linear regression requires categorical features
			\item Linear regression is always accurate
		\end{enumerate}
		
		\vspace{0.1cm}
		\textcolor{myblue}{\textbf{ANSWER: B}}
		
		\vspace{0.1cm}
		\textbf{DETAILED EXPLANATION:}
		\begin{itemize}
			\item Linear regression can predict $< 0$ or $> 1$
			\item No constraint to keep predictions between 0 and 1
			\item 0 and 1 are not probabilities
			\item Logistic regression solves this with sigmoid function
		\end{itemize}
		
		\vspace{0.1cm}
		\textcolor{myred}{\textbf{WHY OTHER OPTIONS ARE WRONG:}}
		\begin{itemize}
			\item \textcolor{myred}{A:} Linear regression is actually simpler
			\item \textcolor{myred}{C:} Linear regression doesn't require categorical features
			\item \textcolor{myred}{D:} Linear regression is NOT accurate for binary outcomes
		\end{itemize}
		
		\textcolor{myred}{\textbf{EXAM TRAP:}} Linear regression can give \textcolor{myblue}{probabilities outside 0-1} range!
	\end{frame}
	
	\begin{frame}{Q29: Logistic Function - Tutorial}
		\fontsize{6.5}{7.5}\selectfont
		\textbf{QUESTION: What is the logistic function also called?}
		
		\vspace{0.1cm}
		\textbf{OPTIONS:}
		\begin{enumerate}
			\item ReLU function
			\item Sigmoid function
			\item Tanh function
			\item Linear function
		\end{enumerate}
		
		\vspace{0.1cm}
		\textcolor{myblue}{\textbf{ANSWER: B}}
		
		\vspace{0.1cm}
		\textbf{DETAILED EXPLANATION:}
		\begin{itemize}
			\item Logistic function = Sigmoid function
			\item Bounded between 0 and 1
			\item S-shaped curve
			\item Formula: $F(y) = \frac{1}{1 + e^{-y}}$
		\end{itemize}
		
		\vspace{0.1cm}
		\textcolor{myred}{\textbf{WHY OTHER OPTIONS ARE WRONG:}}
		\begin{itemize}
			\item \textcolor{myred}{A:} ReLU is a different activation function
			\item \textcolor{myred}{C:} Tanh ranges from -1 to 1
			\item \textcolor{myred}{D:} Linear function is unbounded
		\end{itemize}
		
		\textcolor{myred}{\textbf{EXAM TRAP:}} Logistic = \textcolor{myblue}{Sigmoid} function!
	\end{frame}
	
	\begin{frame}{Q30: Logistic Regression Equation - Tutorial}
		\fontsize{6.5}{7.5}\selectfont
		\textbf{QUESTION: What is the logistic regression model formula for $P(y_i = 1)$?}
		
		\vspace{0.1cm}
		\textbf{OPTIONS:}
		\begin{enumerate}
			\item $P_i = \frac{1}{1 + e^{-(\beta_0 + \beta_1 x_{1i} + ... + \beta_m x_{mi})}}$
			\item $P_i = \beta_0 + \beta_1 x_{1i} + ... + \beta_m x_{mi}$
			\item $P_i = e^{(\beta_0 + \beta_1 x_{1i} + ... + \beta_m x_{mi})}$
			\item $P_i = \frac{1}{1 + e^{(\beta_0 + \beta_1 x_{1i} + ... + \beta_m x_{mi})}}$
		\end{enumerate}
		
		\vspace{0.1cm}
		\textcolor{myblue}{\textbf{ANSWER: A}}
		
		\vspace{0.1cm}
		\textbf{DETAILED EXPLANATION:}
		\begin{itemize}
			\item Logistic function applied to linear combination
			\item Bounds probability between 0 and 1
			\item $\beta_0 + \beta_1 x_{1i} + ... + \beta_m x_{mi}$ is the linear predictor
			\item e is Euler's number ($\approx$ 2.718)
		\end{itemize}
		
		\vspace{0.1cm}
		\textcolor{myred}{\textbf{WHY OTHER OPTIONS ARE WRONG:}}
		\begin{itemize}
			\item \textcolor{myred}{B:} This is linear regression, not logistic
			\item \textcolor{myred}{C:} This is exponential, not logistic
			\item \textcolor{myred}{D:} Wrong sign in the exponent (should be negative)
		\end{itemize}
		
		\textcolor{myred}{\textbf{EXAM TRAP:}} Logistic regression = \textcolor{myblue}{1/(1 + e$^{-linear\_predictor}$)}!
	\end{frame}
	
	\begin{frame}{Q31: Logistic Regression Output - Tutorial}
		\fontsize{6.5}{7.5}\selectfont
		\textbf{QUESTION: What does logistic regression output?}
		
		\vspace{0.1cm}
		\textbf{OPTIONS:}
		\begin{enumerate}
			\item A continuous value with no bounds
			\item A probability between 0 and 1
			\item Either 0 or 1 only
			\item A negative number
		\end{enumerate}
		
		\vspace{0.1cm}
		\textcolor{myblue}{\textbf{ANSWER: B}}
		
		\vspace{0.1cm}
		\textbf{DETAILED EXPLANATION:}
		\begin{itemize}
			\item Outputs probability $P(y_i = 1)$
			\item Always between 0 and 1
			\item Can be thresholded to get binary prediction
			\item Typically threshold = 0.5
		\end{itemize}
		
		\vspace{0.1cm}
		\textcolor{myred}{\textbf{WHY OTHER OPTIONS ARE WRONG:}}
		\begin{itemize}
			\item \textcolor{myred}{A:} Linear regression gives unbounded output
			\item \textcolor{myred}{C:} Output is probability, thresholding gives 0/1
			\item \textcolor{myred}{D:} Probabilities are never negative
		\end{itemize}
		
		\textcolor{myred}{\textbf{EXAM TRAP:}} Logistic regression outputs \textcolor{myblue}{probabilities between 0 and 1}!
	\end{frame}
	
	\begin{frame}{Q32: Logistic Regression Interpretation - Tutorial}
		\fontsize{6.5}{7.5}\selectfont
		\textbf{QUESTION: How are logistic regression coefficients interpreted?}
		
		\vspace{0.1cm}
		\textbf{OPTIONS:}
		\begin{enumerate}
			\item As linear effects on the probability
			\item As linear effects on the log-odds
			\item As linear effects on the target directly
			\item Coefficients cannot be interpreted
		\end{enumerate}
		
		\vspace{0.1cm}
		\textcolor{myblue}{\textbf{ANSWER: B}}
		
		\vspace{0.1cm}
		\textbf{DETAILED EXPLANATION:}
		\begin{itemize}
			\item Logistic transformation is nonlinear
			\item Coefficients affect log-odds linearly
			\item Log-odds = $\log(\frac{P}{1-P})$
			\item Signs and significance are still informative
		\end{itemize}
		
		\vspace{0.1cm}
		\textcolor{myred}{\textbf{WHY OTHER OPTIONS ARE WRONG:}}
		\begin{itemize}
			\item \textcolor{myred}{A:} Effect on probability is nonlinear
			\item \textcolor{myred}{C:} Target is binary, not continuous
			\item \textcolor{myred}{D:} Coefficients CAN be interpreted (log-odds)
		\end{itemize}
		
		\textcolor{myred}{\textbf{EXAM TRAP:}} Logistic coefficients affect \textcolor{myblue}{log-odds}, not probability directly!
	\end{frame}
	
	\begin{frame}{Q33: LendingClub Example - Tutorial}
		\fontsize{6.5}{7.5}\selectfont
		\textbf{QUESTION: In the LendingClub example, which factors significantly affect default probability?}
		
		\vspace{0.1cm}
		\textbf{OPTIONS:}
		\begin{enumerate}
			\item Amount and Income
			\item Term and Interest Rate
			\item Homeowner and Delinquent
			\item Bankruptcies and Employment History
		\end{enumerate}
		
		\vspace{0.1cm}
		\textcolor{myblue}{\textbf{ANSWER: B}}
		
		\vspace{0.1cm}
		\textbf{DETAILED EXPLANATION:}
		\begin{itemize}
			\item Term: longer loans increase default probability
			\item Interest rate: higher rates increase default probability
			\item Both significant at 5% level
			\item Mortgage has negative effect (significant)
		\end{itemize}
		
		\vspace{0.1cm}
		\textcolor{myred}{\textbf{WHY OTHER OPTIONS ARE WRONG:}}
		\begin{itemize}
			\item \textcolor{myred}{A:} Amount and Income are not significant
			\item \textcolor{myred}{C:} Homeowner and Delinquent are not significant
			\item \textcolor{myred}{D:} Bankruptcies and Employment are not significant
		\end{itemize}
		
		\textcolor{myred}{\textbf{EXAM TRAP:}} Term and Interest Rate = \textcolor{myblue}{significant predictors} of default!
	\end{frame}
	
	\begin{frame}{Q34: Mortgage Effect in LendingClub - Tutorial}
		\fontsize{6.5}{7.5}\selectfont
		\textbf{QUESTION: In the LendingClub example, what does a negative coefficient on "Mortgage" mean?}
		
		\vspace{0.1cm}
		\textbf{OPTIONS:}
		\begin{enumerate}
			\item Having a mortgage increases default probability
			\item Having a mortgage decreases default probability
			\item Mortgage has no effect on default
			\item Mortgage is the most important factor
		\end{enumerate}
		
		\vspace{0.1cm}
		\textcolor{myblue}{\textbf{ANSWER: B}}
		
		\vspace{0.1cm}
		\textbf{DETAILED EXPLANATION:}
		\begin{itemize}
			\item Negative coefficient = lower default probability
			\item Mortgage holders are less likely to default
			\item Coefficient is significant at 5% level
			\item May indicate financial stability
		\end{itemize}
		
		\vspace{0.1cm}
		\textcolor{myred}{\textbf{WHY OTHER OPTIONS ARE WRONG:}}
		\begin{itemize}
			\item \textcolor{myred}{A:} Negative means decrease, not increase
			\item \textcolor{myred}{C:} It is significant, so it has an effect
			\item \textcolor{myred}{D:} Other factors may be more important
		\end{itemize}
		
		\textcolor{myred}{\textbf{EXAM TRAP:}} Negative coefficient = \textcolor{myblue}{decrease in default probability}!
	\end{frame}
	
	\begin{frame}{Q35: Statistical Significance - Tutorial}
		\fontsize{6.5}{7.5}\selectfont
		\textbf{QUESTION: What does a coefficient with  denote in regression output?}
		
		\vspace{0.1cm}
		\textbf{OPTIONS:}
		\begin{enumerate}
			\item The coefficient is zero
			\item The coefficient is significant at 1% level
			\item The coefficient is significant at 5% level
			\item The coefficient is not significant
		\end{enumerate}
		
		\vspace{0.1cm}
		\textcolor{myblue}{\textbf{ANSWER: C}}
		
		\vspace{0.1cm}
		\textbf{DETAILED EXPLANATION:}
		\begin{itemize}
			\item  significant at 5\% level
			\item   significant at 1\% level
			\item  significant at 10\% level
			\item Indicates probability that coefficient is zero is low
		\end{itemize}
		
		\vspace{0.1cm}
		\textcolor{myred}{\textbf{WHY OTHER OPTIONS ARE WRONG:}}
		\begin{itemize}
			\item \textcolor{myred}{A:}two  means not zero (significant)
			\item \textcolor{myred}{B:} 3 is 1\% significance
			\item \textcolor{myred}{D:} 2 means it IS significant
		\end{itemize}
		
		\textcolor{myred}{\textbf{EXAM TRAP:}} 2 = \textcolor{myblue}{5\% significance};  3 tar = \textcolor{myblue}{1\% significance}!
	\end{frame}
	
	\begin{frame}{Q36: Use Case - Global Financial Company - Tutorial}
		\fontsize{6.5}{7.5}\selectfont
		\textbf{QUESTION: What was the outcome of the Global Financial Company's logistic regression application?}
		
		\vspace{0.1cm}
		\textbf{OPTIONS:}
		\begin{enumerate}
			\item 5\% increase in business
			\item 9.8\% increase in business
			\item 15\% increase in business
			\item No change in business
		\end{enumerate}
		
		\vspace{0.1cm}
		\textcolor{myblue}{\textbf{ANSWER: B}}
		
		\vspace{0.1cm}
		\textbf{DETAILED EXPLANATION:}
		\begin{itemize}
			\item Used logistic regression to refine lending offers
			\item Predicted likelihood of clients accepting loan proposals
			\item Achieved 9.8\% increase in business
			\item Enabled more targeted and effective offers
		\end{itemize}
		
		\vspace{0.1cm}
		\textcolor{myred}{\textbf{WHY OTHER OPTIONS ARE WRONG:}}
		\begin{itemize}
			\item \textcolor{myred}{A:} Lower than actual increase
			\item \textcolor{myred}{C:} Higher than actual increase
			\item \textcolor{myred}{D:} There was a significant increase
		\end{itemize}
		
		\textcolor{myred}{\textbf{EXAM TRAP:}} Global Financial Company achieved \textcolor{myblue}{9.8\% business increase}!
	\end{frame}
	
	\begin{frame}{Q37: Multinomial Logit Definition - Tutorial}
		\fontsize{6.5}{7.5}\selectfont
		\textbf{QUESTION: What is multinomial logit used for?}
		
		\vspace{0.1cm}
		\textbf{OPTIONS:}
		\begin{enumerate}
			\item Binary classification
			\item Multi-class classification with unordered categories
			\item Multi-class classification with ordered categories
			\item Continuous prediction
		\end{enumerate}
		
		\vspace{0.1cm}
		\textcolor{myblue}{\textbf{ANSWER: B}}
		
		\vspace{0.1cm}
		\textbf{DETAILED EXPLANATION:}
		\begin{itemize}
			\item Extends logistic regression to >2 categories
			\item Categories are nominal (no ordering)
			\item One category is the baseline/reference
			\item Estimates log-odds relative to baseline
		\end{itemize}
		
		\vspace{0.1cm}
		\textcolor{myred}{\textbf{WHY OTHER OPTIONS ARE WRONG:}}
		\begin{itemize}
			\item \textcolor{myred}{A:} Binary is standard logistic regression
			\item \textcolor{myred}{C:} Ordered logit is for ordered categories
			\item \textcolor{myred}{D:} Regression is for continuous prediction
		\end{itemize}
		
		\textcolor{myred}{\textbf{EXAM TRAP:}} Multinomial logit = \textcolor{myblue}{unordered multi-class classification}!
	\end{frame}
	
	\begin{frame}{Q38: Discrete Choice Model Example - Tutorial}
		\fontsize{6.5}{7.5}\selectfont
		\textbf{QUESTION: Which is an example of a discrete choice problem?}
		
		\vspace{0.1cm}
		\textbf{OPTIONS:}
		\begin{enumerate}
			\item Predicting house prices
			\item Commuter choosing between car, bus, bicycle, walking
			\item Predicting stock returns
			\item Estimating company revenue
		\end{enumerate}
		
		\vspace{0.1cm}
		\textcolor{myblue}{\textbf{ANSWER: B}}
		
		\vspace{0.1cm}
		\textbf{DETAILED EXPLANATION:}
		\begin{itemize}
			\item Multiple unordered categories
			\item No inherent ranking between options
			\item Features: age, income, distance, fitness
			\item Multinomial logit is appropriate
		\end{itemize}
		
		\vspace{0.1cm}
		\textcolor{myred}{\textbf{WHY OTHER OPTIONS ARE WRONG:}}
		\begin{itemize}
			\item \textcolor{myred}{A:} House price prediction is regression
			\item \textcolor{myred}{C:} Stock return prediction is regression
			\item \textcolor{myred}{D:} Revenue estimation is regression
		\end{itemize}
		
		\textcolor{myred}{\textbf{EXAM TRAP:}} Discrete choice = \textcolor{myblue}{choosing among unordered categories}!
	\end{frame}
	
	\begin{frame}{Q39: Ordered Logit Definition - Tutorial}
		\fontsize{6.5}{7.5}\selectfont
		\textbf{QUESTION: What is ordered logit used for?}
		
		\vspace{0.1cm}
		\textbf{OPTIONS:}
		\begin{enumerate}
			\item Binary classification
			\item Multi-class with ordered categories
			\item Unordered multi-class classification
			\item Continuous prediction
		\end{enumerate}
		
		\vspace{0.1cm}
		\textcolor{myblue}{\textbf{ANSWER: B}}
		
		\vspace{0.1cm}
		\textbf{DETAILED EXPLANATION:}
		\begin{itemize}
			\item Categories have natural ordering
			\item Examples: Credit ratings (AAA, AA, A, BBB...)
			\item Risk scores (Low, Medium, High)
			\item Uses threshold/cutoff parameters
		\end{itemize}
		
		\vspace{0.1cm}
		\textcolor{myred}{\textbf{WHY OTHER OPTIONS ARE WRONG:}}
		\begin{itemize}
			\item \textcolor{myred}{A:} Binary is standard logistic regression
			\item \textcolor{myred}{C:} Multinomial logit is for unordered
			\item \textcolor{myred}{D:} Regression is for continuous prediction
		\end{itemize}
		
		\textcolor{myred}{\textbf{EXAM TRAP:}} Ordered logit = \textcolor{myblue}{ordered multi-class classification}!
	\end{frame}
	
	\begin{frame}{Q40: Ordinal Data Examples - Tutorial}
		\fontsize{6.5}{7.5}\selectfont
		\textbf{QUESTION: Which is an example of ordinal data?}
		
		\vspace{0.1cm}
		\textbf{OPTIONS:}
		\begin{enumerate}
			\item Loan default (Yes/No)
			\item Credit ratings (AAA, AA, A, BBB)
			\item Commuter transport mode choice
			\item House prices
		\end{enumerate}
		
		\vspace{0.1cm}
		\textcolor{myblue}{\textbf{ANSWER: B}}
		
		\vspace{0.1cm}
		\textbf{DETAILED EXPLANATION:}
		\begin{itemize}
			\item Ordinal data has natural ordering
			\item AAA is better than AA, AA better than A
			\item Ordered logit is appropriate
			\item Other examples: survey responses, risk scores
		\end{itemize}
		
		\vspace{0.1cm}
		\textcolor{myred}{\textbf{WHY OTHER OPTIONS ARE WRONG:}}
		\begin{itemize}
			\item \textcolor{myred}{A:} Binary (no ordering)
			\item \textcolor{myred}{C:} Nominal (no ordering)
			\item \textcolor{myred}{D:} Continuous (not categorical)
		\end{itemize}
		
		\textcolor{myred}{\textbf{EXAM TRAP:}} Credit ratings = \textcolor{myblue}{ordinal data} with natural ordering!
	\end{frame}
	
	\begin{frame}{Q41: Log-Odds Interpretation - Tutorial}
		\fontsize{6.5}{7.5}\selectfont
		\textbf{QUESTION: What does a log-odds of zero indicate?}
		
		\vspace{0.1cm}
		\textbf{OPTIONS:}
		\begin{enumerate}
			\item The category is infinitely more likely
			\item The category is equally likely as the baseline
			\item The category is impossible
			\item The category is certain
		\end{enumerate}
		
		\vspace{0.1cm}
		\textcolor{myblue}{\textbf{ANSWER: B}}
		
		\vspace{0.1cm}
		\textbf{DETAILED EXPLANATION:}
		\begin{itemize}
			\item Log-odds = $\log(\frac{P}{1-P})$
			\item Zero means $\frac{P}{1-P} = 1$
			\item So $P = 0.5$
			\item Equal probability of both outcomes
		\end{itemize}
		
		\vspace{0.1cm}
		\textcolor{myred}{\textbf{WHY OTHER OPTIONS ARE WRONG:}}
		\begin{itemize}
			\item \textcolor{myred}{A:} Positive means more likely
			\item \textcolor{myred}{C:} Zero means equal, not impossible
			\item \textcolor{myred}{D:} Zero means equal, not certain
		\end{itemize}
		
		\textcolor{myred}{\textbf{EXAM TRAP:}} Log-odds of 0 = \textcolor{myblue}{equal probability}!
	\end{frame}
	
	\begin{frame}{Q42: Prior Probability in LDA - Tutorial}
		\fontsize{6.5}{7.5}\selectfont
		\textbf{QUESTION: What is the prior probability in LDA?}
		
		\vspace{0.1cm}
		\textbf{OPTIONS:}
		\begin{enumerate}
			\item Probability of class before seeing data
			\item Probability of class after seeing data
			\item The feature distribution
			\item The error term
		\end{enumerate}
		
		\vspace{0.1cm}
		\textcolor{myblue}{\textbf{ANSWER: A}}
		
		\vspace{0.1cm}
		\textbf{DETAILED EXPLANATION:}
		\begin{itemize}
			\item $\pi_j$ in LDA formula
			\item Unconditional probability of each class
			\item Can be assumed equal
			\item Or estimated from training data frequencies
		\end{itemize}
		
		\vspace{0.1cm}
		\textcolor{myred}{\textbf{WHY OTHER OPTIONS ARE WRONG:}}
		\begin{itemize}
			\item \textcolor{myred}{B:} This is posterior probability
			\item \textcolor{myred}{C:} This is likelihood
			\item \textcolor{myred}{D:} This is different concept
		\end{itemize}
		
		\textcolor{myred}{\textbf{EXAM TRAP:}} Prior = \textcolor{myblue}{probability before seeing data}!
	\end{frame}
	
	\begin{frame}{Q43: Posterior Probability in LDA - Tutorial}
		\fontsize{6.5}{7.5}\selectfont
		\textbf{QUESTION: How is a new observation classified in LDA?}
		
		\vspace{0.1cm}
		\textbf{OPTIONS:}
		\begin{enumerate}
			\item Randomly assigned
			\item Assigned to class with highest discriminant score
			\item Assigned to class with lowest discriminant score
			\item Assigned to the majority class
		\end{enumerate}
		
		\vspace{0.1cm}
		\textcolor{myblue}{\textbf{ANSWER: B}}
		
		\vspace{0.1cm}
		\textbf{DETAILED EXPLANATION:}
		\begin{itemize}
			\item Calculate $\delta_j(x)$ for each class
			\item Choose class with largest $\delta_j(x)$
			\item This maximizes posterior probability
			\item Discriminant score is linear function of x
		\end{itemize}
		
		\vspace{0.1cm}
		\textcolor{myred}{\textbf{WHY OTHER OPTIONS ARE WRONG:}}
		\begin{itemize}
			\item \textcolor{myred}{A:} Classification is deterministic
			\item \textcolor{myred}{C:} Highest score gives best class
			\item \textcolor{myred}{D:} Uses discriminants, not just majority
		\end{itemize}
		
		\textcolor{myred}{\textbf{EXAM TRAP:}} LDA assigns to \textcolor{myblue}{class with highest discriminant score}!
	\end{frame}
	
	\begin{frame}{Q44: LDA Assumptions - Tutorial}
		\fontsize{6.5}{7.5}\selectfont
		\textbf{QUESTION: What are the key assumptions of LDA?}
		
		\vspace{0.1cm}
		\textbf{OPTIONS:}
		\begin{enumerate}
			\item Features are independent and have equal means
			\item Features are normally distributed with common covariance matrix
			\item Features are uniformly distributed
			\item Features are categorical
		\end{enumerate}
		
		\vspace{0.1cm}
		\textcolor{myblue}{\textbf{ANSWER: B}}
		
		\vspace{0.1cm}
		\textbf{DETAILED EXPLANATION:}
		\begin{itemize}
			\item Multivariate normal distribution for each class
			\item Common covariance matrix across classes
			\item Features are continuous
			\item Different mean vectors for each class
		\end{itemize}
		
		\vspace{0.1cm}
		\textcolor{myred}{\textbf{WHY OTHER OPTIONS ARE WRONG:}}
		\begin{itemize}
			\item \textcolor{myred}{A:} Means are different, not equal
			\item \textcolor{myred}{C:} Features are normal, not uniform
			\item \textcolor{myred}{D:} Features are continuous, not categorical
		\end{itemize}
		
		\textcolor{myred}{\textbf{EXAM TRAP:}} LDA = \textcolor{myblue}{normal distributions + common covariance}!
	\end{frame}
	
	\begin{frame}{Q45: LDA Robustness - Tutorial}
		\fontsize{6.5}{7.5}\selectfont
		\textbf{QUESTION: What happens if LDA assumptions are not perfectly met?}
		
		\vspace{0.1cm}
		\textbf{OPTIONS:}
		\begin{enumerate}
			\item LDA cannot be used at all
			\item LDA often remains robust and performs well
			\item LDA always fails completely
			\item LDA requires no assumptions
		\end{enumerate}
		
		\vspace{0.1cm}
		\textcolor{myblue}{\textbf{ANSWER: B}}
		
		\vspace{0.1cm}
		\textbf{DETAILED EXPLANATION:}
		\begin{itemize}
			\item Assumptions are ideal conditions
			\item In practice, often not perfectly met
			\item LDA is quite robust to violations
			\item Often performs well in practice
		\end{itemize}
		
		\vspace{0.1cm}
		\textcolor{myred}{\textbf{WHY OTHER OPTIONS ARE WRONG:}}
		\begin{itemize}
			\item \textcolor{myred}{A:} LDA can still be used
			\item \textcolor{myred}{C:} LDA often works well
			\item \textcolor{myred}{D:} LDA does have assumptions
		\end{itemize}
		
		\textcolor{myred}{\textbf{EXAM TRAP:}} LDA is \textcolor{myblue}{robust} even when assumptions are violated!
	\end{frame}
	
	% ============================================================
	% SECTION 4: LINEAR DISCRIMINANT ANALYSIS - QUESTIONS 46-55
	% ============================================================
	
	\begin{frame}{Q46: LDA Purpose - Tutorial}
		\fontsize{6.5}{7.5}\selectfont
		\textbf{QUESTION: What is the primary purpose of Linear Discriminant Analysis (LDA)?}
		
		\vspace{0.1cm}
		\textbf{OPTIONS:}
		\begin{enumerate}
			\item Regression analysis
			\item Classification with multiple well-separated classes
			\item Feature engineering
			\item Error correction
		\end{enumerate}
		
		\vspace{0.1cm}
		\textcolor{myblue}{\textbf{ANSWER: B}}
		
		\vspace{0.1cm}
		\textbf{DETAILED EXPLANATION:}
		\begin{itemize}
			\item Finds linear combinations that best separate classes
			\item Maximizes between-class separation
			\item Minimizes within-class variance
			\item Alternative to logistic regression for multi-class
		\end{itemize}
		
		\vspace{0.1cm}
		\textcolor{myred}{\textbf{WHY OTHER OPTIONS ARE WRONG:}}
		\begin{itemize}
			\item \textcolor{myred}{A:} LDA is for classification, not regression
			\item \textcolor{myred}{C:} LDA is not for feature engineering
			\item \textcolor{myred}{D:} LDA is not for error correction
		\end{itemize}
		
		\textcolor{myred}{\textbf{EXAM TRAP:}} LDA is for \textcolor{myblue}{classification with multiple classes}!
	\end{frame}
	
	\begin{frame}{Q47: LDA vs Logistic Regression - Tutorial}
		\fontsize{6.5}{7.5}\selectfont
		\textbf{QUESTION: What is a key advantage of LDA over logistic regression?}
		
		\vspace{0.1cm}
		\textbf{OPTIONS:}
		\begin{enumerate}
			\item LDA is always more accurate
			\item LDA handles multi-class problems naturally
			\item LDA requires no assumptions
			\item LDA is faster to compute
		\end{enumerate}
		
		\vspace{0.1cm}
		\textcolor{myblue}{\textbf{ANSWER: B}}
		
		\vspace{0.1cm}
		\textbf{DETAILED EXPLANATION:}
		\begin{itemize}
			\item Logistic regression is naturally binary
			\item LDA extends directly to multiple classes
			\item Multi-class logistic regression requires extensions
			\item LDA gives linear discriminant functions
		\end{itemize}
		
		\vspace{0.1cm}
		\textcolor{myred}{\textbf{WHY OTHER OPTIONS ARE WRONG:}}
		\begin{itemize}
			\item \textcolor{myred}{A:} Neither is always more accurate
			\item \textcolor{myred}{C:} LDA has assumptions (normality, common covariance)
			\item \textcolor{myred}{D:} Speed depends on implementation
		\end{itemize}
		
		\textcolor{myred}{\textbf{EXAM TRAP:}} LDA = \textcolor{myblue}{natural multi-class classification}!
	\end{frame}
	
	\begin{frame}{Q48: LDA Discriminant Function - Tutorial}
		\fontsize{6.5}{7.5}\selectfont
		\textbf{QUESTION: What is the LDA discriminant function formula?}
		
		\vspace{0.1cm}
		\textbf{OPTIONS:}
		\begin{enumerate}
			\item $\delta_j(x) = x^T \Sigma^{-1} \mu_j - \frac{1}{2}\mu_j^T \Sigma^{-1} \mu_j + \log(\pi_j)$
			\item $\delta_j(x) = x^T \Sigma \mu_j$
			\item $\delta_j(x) = \mu_j^T \Sigma^{-1} x$
			\item $\delta_j(x) = \log(\pi_j) + \frac{1}{2}\mu_j^T \Sigma^{-1} \mu_j$
		\end{enumerate}
		
		\vspace{0.1cm}
		\textcolor{myblue}{\textbf{ANSWER: A}}
		
		\vspace{0.1cm}
		\textbf{DETAILED EXPLANATION:}
		\begin{itemize}
			\item $x$: feature vector
			\item $\mu_j$: mean vector for class j
			\item $\Sigma$: common covariance matrix
			\item $\pi_j$: prior probability of class j
			\item First term: x's contribution, linear in x
		\end{itemize}
		
		\vspace{0.1cm}
		\textcolor{myred}{\textbf{WHY OTHER OPTIONS ARE WRONG:}}
		\begin{itemize}
			\item \textcolor{myred}{B:} Missing important terms
			\item \textcolor{myred}{C:} Missing important terms
			\item \textcolor{myred}{D:} Missing the x term
		\end{itemize}
		
		\textcolor{myred}{\textbf{EXAM TRAP:}} LDA discriminant function = \textcolor{myblue}{linear in x with prior adjustment}!
	\end{frame}
	
	\begin{frame}{Q49: LDA Example - Training Data - Tutorial}
		\fontsize{6.5}{7.5}\selectfont
		\textbf{QUESTION: In the LDA example, what are the prior probabilities for default/no-default?}
		
		\vspace{0.1cm}
		\textbf{OPTIONS:}
		\begin{enumerate}
			\item 0.5 and 0.5
			\item 0.7 and 0.3
			\item 0.8 and 0.2
			\item 1.0 and 0.0
		\end{enumerate}
		
		\vspace{0.1cm}
		\textcolor{myblue}{\textbf{ANSWER: B}}
		
		\vspace{0.1cm}
		\textbf{DETAILED EXPLANATION:}
		\begin{itemize}
			\item 7 out of 10 are no default (70%)
			\item 3 out of 10 are default (30%)
			\item Priors based on class frequencies
			\item $\pi_0 = 0.7$ (no default), $\pi_1 = 0.3$ (default)
		\end{itemize}
		
		\vspace{0.1cm}
		\textcolor{myred}{\textbf{WHY OTHER OPTIONS ARE WRONG:}}
		\begin{itemize}
			\item \textcolor{myred}{A:} Not balanced (7 vs 3)
			\item \textcolor{myred}{C:} Not 8 vs 2
			\item \textcolor{myred}{D:} Not all one class\end{itemize}
		
		\textcolor{myred}{\textbf{EXAM TRAP:}} Priors = \textcolor{myblue}{class frequencies in training data} (0.7, 0.3)!
	\end{frame}
	
\begin{frame}{Q50: LDA Example - Class Means - Tutorial}
	\fontsize{6.5}{7.5}\selectfont
	\textbf{QUESTION: In the LDA example, what is the mean balance for non-defaults?}
	
	\vspace{0.1cm}
	\textbf{OPTIONS:}
	\begin{enumerate}
		\item 1.31
		\item -0.56
		\item -0.74
		\item 0.32
	\end{enumerate}
	
	\vspace{0.1cm}
	\textcolor{myblue}{\textbf{ANSWER: B}}
	
	\vspace{0.1cm}
	\textbf{DETAILED EXPLANATION:}
	\begin{itemize}
		\item $\mu_0 = \begin{pmatrix} -0.56 \\ 0.32 \end{pmatrix}$
		\item First element is mean balance for non-defaults
		\item Second element is mean income for non-defaults
	\end{itemize}
	
	\vspace{0.1cm}
	\textcolor{myred}{\textbf{WHY OTHER OPTIONS ARE WRONG:}}
	\begin{itemize}
		\item \textcolor{myred}{A:} This is mean balance for defaults
		\item \textcolor{myred}{C:} This is mean income for defaults
		\item \textcolor{myred}{D:} This is mean income for non-defaults
	\end{itemize}
	
	\textcolor{myred}{\textbf{EXAM TRAP:}} Non-default balance mean = \textcolor{myblue}{-0.56} (standardized)!
\end{frame}

\begin{frame}{Q51: LDA Example - Classification - Tutorial}
	\fontsize{6.5}{7.5}\selectfont
	\textbf{QUESTION: In the LDA example, how is the new borrower classified?}
	
	\vspace{0.1cm}
	\textbf{OPTIONS:}
	\begin{enumerate}
		\item Default
		\item No default
		\item Cannot be classified
		\item Both classes equally likely
	\end{enumerate}
	
	\vspace{0.1cm}
	\textcolor{myblue}{\textbf{ANSWER: B}}
	
	\vspace{0.1cm}
	\textbf{DETAILED EXPLANATION:}
	\begin{itemize}
		\item $\delta_0 = 0.30$ (no default)
		\item $\delta_1 = -3.95$ (default)
		\item Since $\delta_0 > \delta_1$, classify as no default
		\item Discriminant score is much higher for no default
	\end{itemize}
	
	\vspace{0.1cm}
	\textcolor{myred}{\textbf{WHY OTHER OPTIONS ARE WRONG:}}
	\begin{itemize}
		\item \textcolor{myred}{A:} Default score is lower
		\item \textcolor{myred}{C:} We have enough information
		\item \textcolor{myred}{D:} Scores are not equal (0.30 vs -3.95)
	\end{itemize}
	
	\textcolor{myred}{\textbf{EXAM TRAP:}} Classify as \textcolor{myblue}{no default} (higher discriminant score)!
\end{frame}

\begin{frame}{Q52: Covariance Matrix in LDA - Tutorial}
	\fontsize{6.5}{7.5}\selectfont
	\textbf{QUESTION: What does the covariance matrix in LDA measure?}
	
	\vspace{0.1cm}
	\textbf{OPTIONS:}
	\begin{enumerate}
		\item The means of features
		\item The spread and correlation structure of features
		\item The class probabilities
		\item The error term
	\end{enumerate}
	
	\vspace{0.1cm}
	\textcolor{myblue}{\textbf{ANSWER: B}}
	
	\vspace{0.1cm}
	\textbf{DETAILED EXPLANATION:}
	\begin{itemize}
		\item Variances on the diagonal
		\item Covariances off-diagonal
		\item Assumed to be common across classes
		\item $\Sigma$ in the discriminant function
	\end{itemize}
	
	\vspace{0.1cm}
	\textcolor{myred}{\textbf{WHY OTHER OPTIONS ARE WRONG:}}
	\begin{itemize}
		\item \textcolor{myred}{A:} Means are $\mu_j$
		\item \textcolor{myred}{C:} Probabilities are $\pi_j$
		\item \textcolor{myred}{D:} Error is different concept
	\end{itemize}
	
	\textcolor{myred}{\textbf{EXAM TRAP:}} Covariance matrix = \textcolor{myblue}{spread and correlation} of features!
\end{frame}

\begin{frame}{Q53: LDA Example - Standardized Data - Tutorial}
	\fontsize{6.5}{7.5}\selectfont
	\textbf{QUESTION: Why are the features standardized in the LDA example?}
	
	\vspace{0.1cm}
	\textbf{OPTIONS:}
	\begin{enumerate}
		\item To make them categorical
		\item To ensure equal scale for LDA
		\item To remove outliers
		\item To make the data smaller
	\end{enumerate}
	
	\vspace{0.1cm}
	\textcolor{myblue}{\textbf{ANSWER: B}}
	
	\vspace{0.1cm}
	\textbf{DETAILED EXPLANATION:}
	\begin{itemize}
		\item LDA assumes continuous features
		\item Features may have different scales
		\item Standardization gives mean 0, variance 1
		\item Ensures features are comparable
	\end{itemize}
	
	\vspace{0.1cm}
	\textcolor{myred}{\textbf{WHY OTHER OPTIONS ARE WRONG:}}
	\begin{itemize}
		\item \textcolor{myred}{A:} Standardization doesn't make categorical
		\item \textcolor{myred}{C:} Standardization doesn't remove outliers
		\item \textcolor{myred}{D:} Standardization doesn't change data size
	\end{itemize}
	
	\textcolor{myred}{\textbf{EXAM TRAP:}} Standardization = \textcolor{myblue}{equal scales} for LDA!
\end{frame}

\begin{frame}{Q54: LDA and Dimensionality Reduction - Tutorial}
	\fontsize{6.5}{7.5}\selectfont
	\textbf{QUESTION: What is another use of LDA beyond classification?}
	
	\vspace{0.1cm}
	\textbf{OPTIONS:}
	\begin{enumerate}
		\item Feature engineering
		\item Dimensionality reduction
		\item Error detection
		\item Data cleaning
	\end{enumerate}
	
	\vspace{0.1cm}
	\textcolor{myblue}{\textbf{ANSWER: B}}
	
	\vspace{0.1cm}
	\textbf{DETAILED EXPLANATION:}
	\begin{itemize}
		\item Projects data to lower-dimensional space
		\item Maximizes between-class separation
		\item Similar to PCA but supervised
		\item Can reduce to $K-1$ dimensions
	\end{itemize}
	
	\vspace{0.1cm}
	\textcolor{myred}{\textbf{WHY OTHER OPTIONS ARE WRONG:}}
	\begin{itemize}
		\item \textcolor{myred}{A:} LDA uses features, doesn't engineer them
		\item \textcolor{myred}{C:} LDA is not for error detection
		\item \textcolor{myred}{D:} LDA is not for data cleaning
	\end{itemize}
	
	\textcolor{myred}{\textbf{EXAM TRAP:}} LDA can be used for \textcolor{myblue}{dimensionality reduction}!
\end{frame}

\begin{frame}{Q55: LDA vs PCA - Tutorial}
	\fontsize{6.5}{7.5}\selectfont
	\textbf{QUESTION: How does LDA differ from PCA for dimensionality reduction?}
	
	\vspace{0.1cm}
	\textbf{OPTIONS:}
	\begin{enumerate}
		\item LDA is unsupervised; PCA is supervised
		\item LDA is supervised; PCA is unsupervised
		\item Both are supervised
		\item Both are unsupervised
	\end{enumerate}
	
	\vspace{0.1cm}
	\textcolor{myblue}{\textbf{ANSWER: B}}
	
	\vspace{0.1cm}
	\textbf{DETAILED EXPLANATION:}
	\begin{itemize}
		\item LDA uses class labels (supervised)
		\item PCA ignores class labels (unsupervised)
		\item LDA maximizes class separation
		\item PCA maximizes variance
	\end{itemize}
	
	\vspace{0.1cm}
	\textcolor{myred}{\textbf{WHY OTHER OPTIONS ARE WRONG:}}
	\begin{itemize}
		\item \textcolor{myred}{A:} LDA is supervised, not unsupervised
		\item \textcolor{myred}{C:} PCA is unsupervised
		\item \textcolor{myred}{D:} LDA is supervised
	\end{itemize}
	
	\textcolor{myred}{\textbf{EXAM TRAP:}} LDA = \textcolor{myblue}{supervised}; PCA = \textcolor{myblue}{unsupervised}!
\end{frame}


	% ============================================================
	% SECTION 5: PRACTICAL APPLICATIONS - QUESTIONS 56-70
	% ============================================================
	
	\begin{frame}{Q56: Salary Prediction Example - Tutorial}
		\fontsize{6.5}{7.5}\selectfont
		\textbf{QUESTION: In the salary example, what does the intercept of 16.88 mean?}
		
		\vspace{0.1cm}
		\textbf{OPTIONS:}
		\begin{enumerate}
			\item Salary increases by \$16.88 per year
			\item Expected salary with no experience is \$16.88/hr
			\item Maximum possible salary is \$16.88/hr
			\item Salary decreases by \$16.88 per year
		\end{enumerate}
		
		\vspace{0.1cm}
		\textcolor{myblue}{\textbf{ANSWER: B}}
		
		\vspace{0.1cm}
		\textbf{DETAILED EXPLANATION:}
		\begin{itemize}
			\item Intercept beta hat 0 = 16.88\$
			\item Expected salary when experience = 0
			\item "Someone just joining with no experience"
			\item In dollars per hour
		\end{itemize}
		
		\vspace{0.1cm}
		\textcolor{myred}{\textbf{WHY OTHER OPTIONS ARE WRONG:}}
		\begin{itemize}
			\item \textcolor{myred}{A:} This is the slope (1.06)
			\item \textcolor{myred}{C:} Not the maximum possible salary
			\item \textcolor{myred}{D:} Salary increases, doesn't decrease
		\end{itemize}
		
		\textcolor{myred}{\textbf{EXAM TRAP:}} Intercept = \textcolor{myblue}{expected value when feature = 0}!
	\end{frame}
	
	\begin{frame}{Q57: Salary Slope Interpretation - Tutorial}
		\fontsize{6.5}{7.5}\selectfont
		\textbf{QUESTION: In the salary example, what does the slope of 1.06 mean?}
		
		\vspace{0.1cm}
		\textbf{OPTIONS:}
		\begin{enumerate}
			\item Each year of experience increases salary by \$1.06/hr
			\item Each year of experience decreases salary by 1.06/hr
			\item Salary is \$1.06/hr with no experience
			\item Maximum salary is \$1.06/hr
		\end{enumerate}
		
		\vspace{0.1cm}
		\textcolor{myblue}{\textbf{ANSWER: A}}
		
		\vspace{0.1cm}
		\textbf{DETAILED EXPLANATION:}
		\begin{itemize}
			\item Slope beta hat 1 = 1.06\$
			\item For each additional year of experience
			\item Salary increases by \$1.06 per hour
			\item Linear relationship
		\end{itemize}
		
		\vspace{0.1cm}
		\textcolor{myred}{\textbf{WHY OTHER OPTIONS ARE WRONG:}}
		\begin{itemize}
			\item \textcolor{myred}{B:} Slope is positive, so increases
			\item \textcolor{myred}{C:} This is the intercept (16.88)
			\item \textcolor{myred}{D:} Not a maximum
		\end{itemize}
		
		\textcolor{myred}{\textbf{EXAM TRAP:}} Slope = \textcolor{myblue}{change in y for 1-unit change in x}!
	\end{frame}
	
	\begin{frame}{Q58: Quadratic Salary Model - Tutorial}
		\fontsize{6.5}{7.5}\selectfont
		\textbf{QUESTION: Why was a quadratic term added to the salary model?}
		
		\vspace{0.1cm}
		\textbf{OPTIONS:}
		\begin{enumerate}
			\item To make the model linear
			\item To capture a non-linear relationship between experience and salary
			\item To add more features
			\item To reduce model complexity
		\end{enumerate}
		
		\vspace{0.1cm}
		\textcolor{myblue}{\textbf{ANSWER: B}}
		
		\vspace{0.1cm}
		\textbf{DETAILED EXPLANATION:}
		\begin{itemize}
			\item Visual inspection suggested non-linearity
			\item Salary increases quickly then plateaus
			\item Quadratic term allows U-shaped relationship
			\item Better fit to the data
		\end{itemize}
		
		\vspace{0.1cm}
		\textcolor{myred}{\textbf{WHY OTHER OPTIONS ARE WRONG:}}
		\begin{itemize}
			\item \textcolor{myred}{A:} Quadratic makes it non-linear
			\item \textcolor{myred}{C:} It adds a feature but that's not the reason
			\item \textcolor{myred}{D:} It increases complexity
		\end{itemize}
		
		\textcolor{myred}{\textbf{EXAM TRAP:}} Quadratic term captures \textcolor{myblue}{non-linear relationships}!
	\end{frame}
	
	\begin{frame}{Q59: Salary Maximization - Tutorial}
		\fontsize{6.5}{7.5}\selectfont
		\textbf{QUESTION: At what experience level is salary maximized in the quadratic model?}
		
		\vspace{0.1cm}
		\textbf{OPTIONS:}
		\begin{enumerate}
			\item 10.3 years
			\item 17.3 years
			\item 25.0 years
			\item 5.0 years
		\end{enumerate}
		
		\vspace{0.1cm}
		\textcolor{myblue}{\textbf{ANSWER: B}}
		
		\vspace{0.1cm}
		\textbf{DETAILED EXPLANATION:}
		\begin{itemize}
			\item Model: $\hat{y} = 11.82 + 2.77x - 0.08x^2$
			\item Derivative: $\frac{dy}{dx} = 2.77 - 0.16x$
			\item Set to zero: $2.77 - 0.16x = 0$
			\item $x = 2.77 / 0.16 = 17.31$ years
		\end{itemize}
		
		\vspace{0.1cm}
		\textcolor{myred}{\textbf{WHY OTHER OPTIONS ARE WRONG:}}
		\begin{itemize}
			\item \textcolor{myred}{A:} Too low
			\item \textcolor{myred}{C:} Too high
			\item \textcolor{myred}{D:} Too low
		\end{itemize}
		
		\textcolor{myred}{\textbf{EXAM TRAP:}} Maximization occurs at \textcolor{myblue}{17.3 years} of experience!
	\end{frame}
	
	\begin{frame}{Q60: Degree Dummy Interpretation - Tutorial}
		\fontsize{6.5}{7.5}\selectfont
		\textbf{QUESTION: In the salary model with degree dummy, what does 0.92 represent?}
		
		\vspace{0.1cm}
		\textbf{OPTIONS:}
		\begin{enumerate}
			\item Salary increase of \$0.92/hr for having a degree
			\item Salary decrease of \$0.92/hr for having a degree
			\item Experience effect
			\item Intercept
		\end{enumerate}
		
		\vspace{0.1cm}
		\textcolor{myblue}{\textbf{ANSWER: A}}
		
		\vspace{0.1cm}
		\textbf{DETAILED EXPLANATION:}
		\begin{itemize}
			\item Coefficient on degree dummy = 0.92
			\item Degree holders earn \$0.92/hr more
			\item Holding experience constant
			\item "Compared with someone having identical experience but no degree"
		\end{itemize}
		
		\vspace{0.1cm}
		\textcolor{myred}{\textbf{WHY OTHER OPTIONS ARE WRONG:}}
		\begin{itemize}
			\item \textcolor{myred}{B:} Positive means increase
			\item \textcolor{myred}{C:} This is the degree effect
			\item \textcolor{myred}{D:} Intercept is 11.38
		\end{itemize}
		
		\textcolor{myred}{\textbf{EXAM TRAP:}} Dummy coefficient = \textcolor{myblue}{effect of category} compared to baseline!
	\end{frame}
	
	\begin{frame}{Q61: Dummy Variable Definition - Tutorial}
		\fontsize{6.5}{7.5}\selectfont
		\textbf{QUESTION: What is a dummy variable?}
		
		\vspace{0.1cm}
		\textbf{OPTIONS:}
		\begin{enumerate}
			\item A continuous variable
			\item A binary variable (0 or 1) representing categories
			\item A variable with more than two categories
			\item A variable that is always zero
		\end{enumerate}
		
		\vspace{0.1cm}
		\textcolor{myblue}{\textbf{ANSWER: B}}
		
		\vspace{0.1cm}
		\textbf{DETAILED EXPLANATION:}
		\begin{itemize}
			\item Takes values 0 or 1
			\item Represents categorical data
			\item Example: Degree (1 = has degree, 0 = no degree)
			\item Can be used in regression models
		\end{itemize}
		
		\vspace{0.1cm}
		\textcolor{myred}{\textbf{WHY OTHER OPTIONS ARE WRONG:}}
		\begin{itemize}
			\item \textcolor{myred}{A:} Dummy variables are binary
			\item \textcolor{myred}{C:} Multiple categories need multiple dummies
			\item \textcolor{myred}{D:} Dummy variables can be 0 or 1
		\end{itemize}
		
		\textcolor{myred}{\textbf{EXAM TRAP:}} Dummy variables are \textcolor{myblue}{binary (0/1)}!
	\end{frame}
	
	\begin{frame}{Q62: Dummy Variable Trap - Tutorial}
		\fontsize{6.5}{7.5}\selectfont
		\textbf{QUESTION: What is the dummy variable trap?}
		
		\vspace{0.1cm}
		\textbf{OPTIONS:}
		\begin{enumerate}
			\item Including too many dummy variables
			\item Including a set of dummies that sum to one along with an intercept
			\item Not including enough dummy variables
			\item Including only one dummy variable
		\end{enumerate}
		
		\vspace{0.1cm}
		\textcolor{myblue}{\textbf{ANSWER: B}}
		
		\vspace{0.1cm}
		\textbf{DETAILED EXPLANATION:}
		\begin{itemize}
			\item Example: Gender dummies (male, female) + intercept
			\item Dummies sum to 1 for every observation
			\item Perfect multicollinearity
			\item OLS breaks down
		\end{itemize}
		
		\vspace{0.1cm}
		\textcolor{myred}{\textbf{WHY OTHER OPTIONS ARE WRONG:}}
		\begin{itemize}
			\item \textcolor{myred}{A:} Too many is not the trap
			\item \textcolor{myred}{C:} Not enough is not the trap
			\item \textcolor{myred}{D:} One dummy is fine
		\end{itemize}
		
		\textcolor{myred}{\textbf{EXAM TRAP:}} Dummy variable trap = \textcolor{myblue}{perfect multicollinearity} with intercept!
	\end{frame}
	
	\begin{frame}{Q63: Seasonal Dummies - Tutorial}
		\fontsize{6.5}{7.5}\selectfont
		\textbf{QUESTION: What are seasonal dummies used for in time-series regression?}
		
		\vspace{0.1cm}
		\textbf{OPTIONS:}
		\begin{enumerate}
			\item To remove seasonality
			\item To capture patterns that repeat weekly, monthly, or quarterly
			\item To make data stationary
			\item To remove trends
		\end{enumerate}
		
		\vspace{0.1cm}
		\textcolor{myblue}{\textbf{ANSWER: B}}
		
		\vspace{0.1cm}
		\textbf{DETAILED EXPLANATION:}
		\begin{itemize}
			\item Capture repeating patterns
			\item Weekly, monthly, quarterly effects
			\item Example: Retail sales higher in December
			\item Control for seasonal variation
		\end{itemize}
		
		\vspace{0.1cm}
		\textcolor{myred}{\textbf{WHY OTHER OPTIONS ARE WRONG:}}
		\begin{itemize}
			\item \textcolor{myred}{A:} They capture, not remove, seasonality
			\item \textcolor{myred}{C:} This is a different concept
			\item \textcolor{myred}{D:} Trends are different from seasonality
		\end{itemize}
		
		\textcolor{myred}{\textbf{EXAM TRAP:}} Seasonal dummies = \textcolor{myblue}{capture repeating patterns}!
	\end{frame}
	
	\begin{frame}{Q64: LOG Transformations - Tutorial}
		\fontsize{6.5}{7.5}\selectfont
		\textbf{QUESTION: Why might you apply a logarithmic transformation to variables?}
		
		\vspace{0.1cm}
		\textbf{OPTIONS:}
		\begin{enumerate}
			\item To make the model non-linear
			\item To handle different scales or non-linear relationships
			\item To remove all variation
			\item To make variables categorical
		\end{enumerate}
		
		\vspace{0.1cm}
		\textcolor{myblue}{\textbf{ANSWER: B}}
		
		\vspace{0.1cm}
		\textbf{DETAILED EXPLANATION:}
		\begin{itemize}
			\item Handles different scales
			\item Can address heteroskedasticity
			\item Changes interpretation of coefficients
			\item Model remains linear in parameters
		\end{itemize}
		
		\vspace{0.1cm}
		\textcolor{myred}{\textbf{WHY OTHER OPTIONS ARE WRONG:}}
		\begin{itemize}
			\item \textcolor{myred}{A:} Model is still linear in parameters
			\item \textcolor{myred}{C:} Variation remains, but transformed
			\item \textcolor{myred}{D:} Variables remain continuous
		\end{itemize}
		
		\textcolor{myred}{\textbf{EXAM TRAP:}} Log transformations = \textcolor{myblue}{handle different scales and non-linearity}!
	\end{frame}
	
	\begin{frame}{Q65: R-squared Interpretation - Tutorial}
		\fontsize{6.5}{7.5}\selectfont
		\textbf{QUESTION: What does R-squared measure in regression?}
		
		\vspace{0.1cm}
		\textbf{OPTIONS:}
		\begin{enumerate}
			\item The proportion of variance explained by the model
			\item The significance of coefficients
			\item The correlation between features
			\item The error term
		\end{enumerate}
		
		\vspace{0.1cm}
		\textcolor{myblue}{\textbf{ANSWER: A}}
		
		\vspace{0.1cm}
		\textbf{DETAILED EXPLANATION:}
		\begin{itemize}
			\item Ranges from 0 to 1
			\item Higher = better fit
			\item Proportion of total variation explained
			\item $R^2 = 1 - \frac{SSE}{SST}$
		\end{itemize}
		
		\vspace{0.1cm}
		\textcolor{myred}{\textbf{WHY OTHER OPTIONS ARE WRONG:}}
		\begin{itemize}
			\item \textcolor{myred}{B:} Significance is p-values
			\item \textcolor{myred}{C:} Correlation is different
			\item \textcolor{myred}{D:} Error term is separate
		\end{itemize}
		
		\textcolor{myred}{\textbf{EXAM TRAP:}} R-squared = \textcolor{myblue}{variance explained by model}!
	\end{frame}
	
	\begin{frame}{Q66: Adjusted R-squared - Tutorial}
		\fontsize{6.5}{7.5}\selectfont
		\textbf{QUESTION: Why use adjusted R-squared instead of R-squared?}
		
		\vspace{0.1cm}
		\textbf{OPTIONS:}
		\begin{enumerate}
			\item It is always larger than R-squared
			\item It penalizes model complexity (number of features)
			\item It measures significance
			\item It is easier to calculate
		\end{enumerate}
		
		\vspace{0.1cm}
		\textcolor{myblue}{\textbf{ANSWER: B}}
		
		\vspace{0.1cm}
		\textbf{DETAILED EXPLANATION:}
		\begin{itemize}
			\item R-squared always increases with more features
			\item Adjusted R-squared penalizes adding features
			\item Formula: $\bar{R}^2 = 1 - \frac{(1-R^2)(n-1)}{n-k-1}$
			\item Helps prevent overfitting
		\end{itemize}
		
		\vspace{0.1cm}
		\textcolor{myred}{\textbf{WHY OTHER OPTIONS ARE WRONG:}}
		\begin{itemize}
			\item \textcolor{myred}{A:} Adjusted R-squared is usually smaller
			\item \textcolor{myred}{C:} Significance is p-values
			\item \textcolor{myred}{D:} R-squared is easier to calculate
		\end{itemize}
		
		\textcolor{myred}{\textbf{EXAM TRAP:}} Adjusted R-squared = \textcolor{myblue}{penalizes model complexity}!
	\end{frame}
	
	\begin{frame}{Q67: Overfitting Definition - Tutorial}
		\fontsize{6.5}{7.5}\selectfont
		\textbf{QUESTION: What is overfitting in regression?}
		
		\vspace{0.1cm}
		\textbf{OPTIONS:}
		\begin{enumerate}
			\item Model fits training data perfectly
			\item Model captures noise instead of true relationship
			\item Model fits test data better
			\item Model has too few features
		\end{enumerate}
		
		\vspace{0.1cm}
		\textcolor{myblue}{\textbf{ANSWER: B}}
		
		\vspace{0.1cm}
		\textbf{DETAILED EXPLANATION:}
		\begin{itemize}
			\item Too complex model
			\item Fits noise in training data
			\item Poor performance on new data
			\item Caused by too many features or high-order terms
		\end{itemize}
		
		\vspace{0.1cm}
		\textcolor{myred}{\textbf{WHY OTHER OPTIONS ARE WRONG:}}
		\begin{itemize}
			\item \textcolor{myred}{A:} Overfitting fits training perfectly but poorly generalizes
			\item \textcolor{myred}{C:} Overfitting fails on test data
			\item \textcolor{myred}{D:} Overfitting has too many features
		\end{itemize}
		
		\textcolor{myred}{\textbf{EXAM TRAP:}} Overfitting = \textcolor{myblue}{capturing noise, not signal}!
	\end{frame}
	
	\begin{frame}{Q68: Underfitting Definition - Tutorial}
		\fontsize{6.5}{7.5}\selectfont
		\textbf{QUESTION: What is underfitting in regression?}
		
		\vspace{0.1cm}
		\textbf{OPTIONS:}
		\begin{enumerate}
			\item Model is too complex
			\item Model is too simple to capture relationships
			\item Model fits training data perfectly
			\item Model has too many features
		\end{enumerate}
		
		\vspace{0.1cm}
		\textcolor{myblue}{\textbf{ANSWER: B}}
		
		\vspace{0.1cm}
		\textbf{DETAILED EXPLANATION:}
		\begin{itemize}
			\item Model is too simple
			\item Cannot capture true relationship
			\item High bias, low variance
			\item Example: Linear model for quadratic data
		\end{itemize}
		
		\vspace{0.1cm}
		\textcolor{myred}{\textbf{WHY OTHER OPTIONS ARE WRONG:}}
		\begin{itemize}
			\item \textcolor{myred}{A:} This is overfitting
			\item \textcolor{myred}{C:} This is overfitting
			\item \textcolor{myred}{D:} This is overfitting
		\end{itemize}
		
		\textcolor{myred}{\textbf{EXAM TRAP:}} Underfitting = \textcolor{myblue}{model too simple}!
	\end{frame}
	
	\begin{frame}{Q69: Bias-Variance Tradeoff - Tutorial}
		\fontsize{6.5}{7.5}\selectfont
		\textbf{QUESTION: What is the bias-variance tradeoff?}
		
		\vspace{0.1cm}
		\textbf{OPTIONS:}
		\begin{enumerate}
			\item Tradeoff between model complexity and interpretability
			\item Tradeoff between model simplicity (bias) and sensitivity to data (variance)
			\item Tradeoff between training and test error
			\item Tradeoff between features and observations
		\end{enumerate}
		
		\vspace{0.1cm}
		\textcolor{myblue}{\textbf{ANSWER: B}}
		
		\vspace{0.1cm}
		\textbf{DETAILED EXPLANATION:}
		\begin{itemize}
			\item Simple models: high bias, low variance
			\item Complex models: low bias, high variance
			\item Need to balance for optimal performance
			\item Central to machine learning
		\end{itemize}
		
		\vspace{0.1cm}
		\textcolor{myred}{\textbf{WHY OTHER OPTIONS ARE WRONG:}}
		\begin{itemize}
			\item \textcolor{myred}{A:} This is accuracy-interpretability tradeoff
			\item \textcolor{myred}{C:} This is a consequence, not definition
			\item \textcolor{myred}{D:} This is different concept
		\end{itemize}
		
		\textcolor{myred}{\textbf{EXAM TRAP:}} Bias-variance = \textcolor{myblue}{simplicity vs complexity}!
	\end{frame}
	
	\begin{frame}{Q70: Regularization Introduction - Tutorial}
		\fontsize{6.5}{7.5}\selectfont
		\textbf{QUESTION: What is regularization in regression?}
		
		\vspace{0.1cm}
		\textbf{OPTIONS:}
		\begin{enumerate}
			\item Adding penalty terms to reduce overfitting
			\item Removing all features
			\item Increasing model complexity
			\item Ignoring errors
		\end{enumerate}
		
		\vspace{0.1cm}
		\textcolor{myblue}{\textbf{ANSWER: A}}
		
		\vspace{0.1cm}
		\textbf{DETAILED EXPLANATION:}
		\begin{itemize}
			\item Adds penalty on coefficients
			\item Reduces model complexity
			\item Helps prevent overfitting
			\item Examples: LASSO, Ridge (Chapter 7)
		\end{itemize}
		
		\vspace{0.1cm}
		\textcolor{myred}{\textbf{WHY OTHER OPTIONS ARE WRONG:}}
		\begin{itemize}
			\item \textcolor{myred}{B:} Regularization shrinks, doesn't remove all
			\item \textcolor{myred}{C:} Regularization REDUCES complexity
			\item \textcolor{myred}{D:} Regularization penalizes errors
		\end{itemize}
		
		\textcolor{myred}{\textbf{EXAM TRAP:}} Regularization = \textcolor{myblue}{penalty to reduce overfitting}!
	\end{frame}
	
	% ============================================================
	% SECTION 6: ADVANCED TOPICS - QUESTIONS 71-85
	% ============================================================
	
	\begin{frame}{Q71: AIC Definition - Tutorial}
		\fontsize{6.5}{7.5}\selectfont
		\textbf{QUESTION: What does AIC measure in stepwise regression?}
		
		\vspace{0.1cm}
		\textbf{OPTIONS:}
		\begin{enumerate}
			\item Model accuracy
			\item Prediction error penalized for model complexity
			\item Correlation between features
			\item Significance of coefficients
		\end{enumerate}
		
		\vspace{0.1cm}
		\textcolor{myblue}{\textbf{ANSWER: B}}
		
		\vspace{0.1cm}
		\textbf{DETAILED EXPLANATION:}
		\begin{itemize}
			\item AIC = Akaike Information Criterion
			\item Measures prediction error
			\item Penalizes models with more features
			\item Lower AIC indicates better model
		\end{itemize}
		
		\vspace{0.1cm}
		\textcolor{myred}{\textbf{WHY OTHER OPTIONS ARE WRONG:}}
		\begin{itemize}
			\item \textcolor{myred}{A:} AIC is not directly accuracy
			\item \textcolor{myred}{C:} Correlation is separate
			\item \textcolor{myred}{D:} Significance is p-values
		\end{itemize}
		
		\textcolor{myred}{\textbf{EXAM TRAP:}} AIC = \textcolor{myblue}{prediction error + complexity penalty}!
	\end{frame}
	
	\begin{frame}{Q72: Forward Selection Process - Tutorial}
		\fontsize{6.5}{7.5}\selectfont
		\textbf{QUESTION: In forward selection, when does the process stop?}
		
		\vspace{0.1cm}
		\textbf{OPTIONS:}
		\begin{enumerate}
			\item When all features are included
			\item When adding any feature no longer improves the criterion
			\item After exactly 3 steps
			\item When AIC reaches zero
		\end{enumerate}
		
		\vspace{0.1cm}
		\textcolor{myblue}{\textbf{ANSWER: B}}
		
		\vspace{0.1cm}
		\textbf{DETAILED EXPLANATION:}
		\begin{itemize}
			\item Start with intercept only
			\item Add features one by one
			\item Stop when improvement stops
			\item Or when AIC starts increasing
		\end{itemize}
		
		\vspace{0.1cm}
		\textcolor{myred}{\textbf{WHY OTHER OPTIONS ARE WRONG:}}
		\begin{itemize}
			\item \textcolor{myred}{A:} May stop before all features
			\item \textcolor{myred}{C:} Not a fixed number of steps
			\item \textcolor{myred}{D:} AIC is never zero in practice
		\end{itemize}
		
		\textcolor{myred}{\textbf{EXAM TRAP:}} Forward selection stops when \textcolor{myblue}{no improvement in criterion}!
	\end{frame}
	
	\begin{frame}{Q73: Backward Selection Process - Tutorial}
		\fontsize{6.5}{7.5}\selectfont
		\textbf{QUESTION: In backward selection, when does the process stop?}
		
		\vspace{0.1cm}
		\textbf{OPTIONS:}
		\begin{enumerate}
			\item When all features are removed
			\item When removing any feature no longer improves the criterion
			\item After exactly 3 steps
			\item When AIC reaches zero
		\end{enumerate}
		
		\vspace{0.1cm}
		\textcolor{myblue}{\textbf{ANSWER: B}}
		
		\vspace{0.1cm}
		\textbf{DETAILED EXPLANATION:}
		\begin{itemize}
			\item Start with all features
			\item Remove features one by one
			\item Stop when removal worsens criterion
			\item Or when AIC starts increasing
		\end{itemize}
		
		\vspace{0.1cm}
		\textcolor{myred}{\textbf{WHY OTHER OPTIONS ARE WRONG:}}
		\begin{itemize}
			\item \textcolor{myred}{A:} May stop before all removed
			\item \textcolor{myred}{C:} Not a fixed number of steps
			\item \textcolor{myred}{D:} AIC is never zero in practice
		\end{itemize}
		
		\textcolor{myred}{\textbf{EXAM TRAP:}} Backward selection stops when \textcolor{myblue}{removal no longer improves}!
	\end{frame}
	
	\begin{frame}{Q74: Stepwise Regression Limitation - Tutorial}
		\fontsize{6.5}{7.5}\selectfont
		\textbf{QUESTION: What is a limitation of stepwise regression?}
		
		\vspace{0.1cm}
		\textbf{OPTIONS:}
		\begin{enumerate}
			\item It always finds the best model
			\item It may not find the absolute best model and explores limited paths
			\item It is computationally impossible
			\item It requires no domain knowledge
		\end{enumerate}
		
		\vspace{0.1cm}
		\textcolor{myblue}{\textbf{ANSWER: B}}
		
		\vspace{0.1cm}
		\textbf{DETAILED EXPLANATION:}
		\begin{itemize}
			\item Heuristic approach
			\item May miss optimal combinations
			\item Forward and backward may give different results
			\item Not guaranteed to find absolute best subset
		\end{itemize}
		
		\vspace{0.1cm}
		\textcolor{myred}{\textbf{WHY OTHER OPTIONS ARE WRONG:}}
		\begin{itemize}
			\item \textcolor{myred}{A:} It may NOT find the best model
			\item \textcolor{myred}{C:} It is computationally feasible
			\item \textcolor{myred}{D:} Domain knowledge is still needed
		\end{itemize}
		
		\textcolor{myred}{\textbf{EXAM TRAP:}} Stepwise is \textcolor{myblue}{heuristic}, not guaranteed optimal!
	\end{frame}
	
	\begin{frame}{Q75: Best Subset Selection - Tutorial}
		\fontsize{6.5}{7.5}\selectfont
		\textbf{QUESTION: Why is best subset selection impractical for many features?}
		
		\vspace{0.1cm}
		\textbf{OPTIONS:}
		\begin{enumerate}
			\item Because it's too simple
			\item Because there are $2^m$ possible models to evaluate
			\item Because it uses no criteria
			\item Because it requires no data
		\end{enumerate}
		
		\vspace{0.1cm}
		\textcolor{myblue}{\textbf{ANSWER: B}}
		
		\vspace{0.1cm}
		\textbf{DETAILED EXPLANATION:}
		\begin{itemize}
			\item For m features, $2^m$ possible models
			\item With m=20, over 1 million models
			\item Computationally expensive
			\item Stepwise is a practical alternative
		\end{itemize}
		
		\vspace{0.1cm}
		\textcolor{myred}{\textbf{WHY OTHER OPTIONS ARE WRONG:}}
		\begin{itemize}
			\item \textcolor{myred}{A:} It is computationally complex
			\item \textcolor{myred}{C:} It uses criteria like AIC
			\item \textcolor{myred}{D:} It requires data
		\end{itemize}
		
		\textcolor{myred}{\textbf{EXAM TRAP:}} Best subset has \textcolor{myblue}{$2^m$ models} to evaluate!
	\end{frame}
	
	\begin{frame}{Q76: Binary Outcomes in Risk - Tutorial}
		\fontsize{6.5}{7.5}\selectfont
		\textbf{QUESTION: Which is NOT a binary outcome in risk management?}
		
		\vspace{0.1cm}
		\textbf{OPTIONS:}
		\begin{enumerate}
			\item Loan defaults
			\item Fraud detection
			\item Employee salary
			\item Equipment failure
		\end{enumerate}
		
		\vspace{0.1cm}
		\textcolor{myblue}{\textbf{ANSWER: C}}
		
		\vspace{0.1cm}
		\textbf{DETAILED EXPLANATION:}
		\begin{itemize}
			\item Employee salary is continuous
			\item Default, fraud, failure are binary
			\item Binary outcomes are Yes/No
			\item Salary is a numerical value
		\end{itemize}
		
		\vspace{0.1cm}
		\textcolor{myred}{\textbf{WHY OTHER OPTIONS ARE WRONG:}}
		\begin{itemize}
			\item \textcolor{myred}{A:} Default = Yes/No
			\item \textcolor{myred}{B:} Fraud = Yes/No
			\item \textcolor{myred}{D:} Failure = Yes/No
		\end{itemize}
		
		\textcolor{myred}{\textbf{EXAM TRAP:}} Salary is \textcolor{myblue}{continuous}, not binary!
	\end{frame}
	
	\begin{frame}{Q77: Logistic Regression Applications - Tutorial}
		\fontsize{6.5}{7.5}\selectfont
		\textbf{QUESTION: What is a common application of logistic regression?}
		
		\vspace{0.1cm}
		\textbf{OPTIONS:}
		\begin{enumerate}
			\item Predicting house prices
			\item Credit risk assessment (default prediction)
			\item Stock return forecasting
			\item Estimating company revenue
		\end{enumerate}
		
		\vspace{0.1cm}
		\textcolor{myblue}{\textbf{ANSWER: B}}
		
		\vspace{0.1cm}
		\textbf{DETAILED EXPLANATION:}
		\begin{itemize}
			\item Logistic regression for binary outcomes
			\item Default prediction is binary (default/no default)
			\item Credit risk is common application
			\item Example: LendingClub case
		\end{itemize}
		
		\vspace{0.1cm}
		\textcolor{myred}{\textbf{WHY OTHER OPTIONS ARE WRONG:}}
		\begin{itemize}
			\item \textcolor{myred}{A:} House prices are continuous (regression)
			\item \textcolor{myred}{C:} Stock returns are continuous (regression)
			\item \textcolor{myred}{D:} Revenue is continuous (regression)
		\end{itemize}
		
		\textcolor{myred}{\textbf{EXAM TRAP:}} Logistic regression = \textcolor{myblue}{binary classification} like default prediction!
	\end{frame}
	
	\begin{frame}{Q78: Nominal Data - Tutorial}
		\fontsize{6.5}{7.5}\selectfont
		\textbf{QUESTION: What is nominal data?}
		
		\vspace{0.1cm}
		\textbf{OPTIONS:}
		\begin{enumerate}
			\item Data with natural ordering
			\item Data with categories but no ordering
			\item Continuous numerical data
			\item Binary data
		\end{enumerate}
		
		\vspace{0.1cm}
		\textcolor{myblue}{\textbf{ANSWER: B}}
		
		\vspace{0.1cm}
		\textbf{DETAILED EXPLANATION:}
		\begin{itemize}
			\item Categories have no inherent ranking
			\item Example: Transport modes (car, bus, bicycle)
			\item Multinomial logit is appropriate
			\item Different from ordinal data
		\end{itemize}
		
		\vspace{0.1cm}
		\textcolor{myred}{\textbf{WHY OTHER OPTIONS ARE WRONG:}}
		\begin{itemize}
			\item \textcolor{myred}{A:} This is ordinal data
			\item \textcolor{myred}{C:} This is continuous data
			\item \textcolor{myred}{D:} Binary is a special case
		\end{itemize}
		
		\textcolor{myred}{\textbf{EXAM TRAP:}} Nominal = \textcolor{myblue}{categories without ordering}!
	\end{frame}
	
	\begin{frame}{Q79: Ordinal Data Examples - Tutorial}
		\fontsize{6.5}{7.5}\selectfont
		\textbf{QUESTION: Which is an example of ordinal data?}
		
		\vspace{0.1cm}
		\textbf{OPTIONS:}
		\begin{enumerate}
			\item Customer satisfaction (Very satisfied, Satisfied, Neutral, Unsatisfied)
			\item Color of cars
			\item City of residence
			\item Employee ID number
		\end{enumerate}
		
		\vspace{0.1cm}
		\textcolor{myblue}{\textbf{ANSWER: A}}
		
		\vspace{0.1cm}
		\textbf{DETAILED EXPLANATION:}
		\begin{itemize}
			\item Has natural ordering
			\item Very satisfied > Satisfied > Neutral > Unsatisfied
			\item Ordered logit is appropriate
			\item Other examples: credit ratings, risk scores
		\end{itemize}
		
		\vspace{0.1cm}
		\textcolor{myred}{\textbf{WHY OTHER OPTIONS ARE WRONG:}}
		\begin{itemize}
			\item \textcolor{myred}{B:} Colors have no ordering (nominal)
			\item \textcolor{myred}{C:} Cities have no ordering (nominal)
			\item \textcolor{myred}{D:} IDs have no ordering (nominal)
		\end{itemize}
		
		\textcolor{myred}{\textbf{EXAM TRAP:}} Ordinal = \textcolor{myblue}{ordered categories}!
	\end{frame}
	
	\begin{frame}{Q80: LDA Prior Probabilities - Tutorial}
		\fontsize{6.5}{7.5}\selectfont
		\textbf{QUESTION: How are prior probabilities estimated in LDA?}
		
		\vspace{0.1cm}
		\textbf{OPTIONS:}
		\begin{enumerate}
			\item Always set to 0.5
			\item Estimated from class frequencies in training data
			\item Randomly assigned
			\item Set to 1
		\end{enumerate}
		
		\vspace{0.1cm}
		\textcolor{myblue}{\textbf{ANSWER: B}}
		
		\vspace{0.1cm}
		\textbf{DETAILED EXPLANATION:}
		\begin{itemize}
			\item $\pi_j = \frac{n_j}{N}$
			\item $n_j$: number of training observations in class j
			\item $N$: total training observations
			\item Can also be assumed equal if desired
		\end{itemize}
		
		\vspace{0.1cm}
		\textcolor{myred}{\textbf{WHY OTHER OPTIONS ARE WRONG:}}
		\begin{itemize}
			\item \textcolor{myred}{A:} Not always 0.5
			\item \textcolor{myred}{C:} Not random
			\item \textcolor{myred}{D:} Not always 1
		\end{itemize}
		
		\textcolor{myred}{\textbf{EXAM TRAP:}} Priors = \textcolor{myblue}{class frequencies} in training data!
	\end{frame}
	
	\begin{frame}{Q81: LDA Discriminant Scores - Tutorial}
		\fontsize{6.5}{7.5}\selectfont
		\textbf{QUESTION: In LDA, what determines the classification?}
		
		\vspace{0.1cm}
		\textbf{OPTIONS:}
		\begin{enumerate}
			\item The lowest discriminant score
			\item The highest discriminant score
			\item The mean of all scores
			\item Random selection
		\end{enumerate}
		
		\vspace{0.1cm}
		\textcolor{myblue}{\textbf{ANSWER: B}}
		
		\vspace{0.1cm}
		\textbf{DETAILED EXPLANATION:}
		\begin{itemize}
			\item Calculate $\delta_j(x)$ for each class
			\item Choose class with maximum $\delta_j(x)$
			\item This maximizes posterior probability
			\item Linear in features
		\end{itemize}
		
		\vspace{0.1cm}
		\textcolor{myred}{\textbf{WHY OTHER OPTIONS ARE WRONG:}}
		\begin{itemize}
			\item \textcolor{myred}{A:} Lowest score would be worst
			\item \textcolor{myred}{C:} Mean doesn't determine classification
			\item \textcolor{myred}{D:} Classification is deterministic
		\end{itemize}
		
		\textcolor{myred}{\textbf{EXAM TRAP:}} LDA chooses \textcolor{myblue}{highest discriminant score}!
	\end{frame}
	
	\begin{frame}{Q82: Multinomial Logit vs Ordered Logit - Tutorial}
		\fontsize{6.5}{7.5}\selectfont
		\textbf{QUESTION: When would you use ordered logit instead of multinomial logit?}
		
		\vspace{0.1cm}
		\textbf{OPTIONS:}
		\begin{enumerate}
			\item When categories have natural ordering
			\item When categories are unordered
			\item When predicting continuous values
			\item When there are only two categories
		\end{enumerate}
		
		\vspace{0.1cm}
		\textcolor{myblue}{\textbf{ANSWER: A}}
		
		\vspace{0.1cm}
		\textbf{DETAILED EXPLANATION:}
		\begin{itemize}
			\item Ordered logit for ordinal data
			\item Multinomial logit for nominal data
			\item Ordered logit preserves ordering information
			\item Example: Credit ratings (AAA, AA, A, BBB)
		\end{itemize}
		
		\vspace{0.1cm}
		\textcolor{myred}{\textbf{WHY OTHER OPTIONS ARE WRONG:}}
		\begin{itemize}
			\item \textcolor{myred}{B:} Multinomial is for unordered
			\item \textcolor{myred}{C:} Neither is for continuous
			\item \textcolor{myred}{D:} Binary logistic for two categories
		\end{itemize}
		
		\textcolor{myred}{\textbf{EXAM TRAP:}} Ordered logit = \textcolor{myblue}{ordered categories}!
	\end{frame}
	
\begin{frame}{Q83: LDA Example Calculation - Tutorial}
	\fontsize{6.5}{7.5}\selectfont
	\textbf{QUESTION: In the LDA example, what is the inverse of the covariance matrix?}
	
	\vspace{0.1cm}
	\textbf{OPTIONS:}
	\begin{enumerate}
		\item $\begin{pmatrix} 1 & 0.58 \\ 0.58 & 1 \end{pmatrix}$
		\item $\begin{pmatrix} 1.52 & 0.88 \\ 0.88 & 1.52 \end{pmatrix}$
		\item $\begin{pmatrix} -1.52 & -0.88 \\ -0.88 & -1.52 \end{pmatrix}$
		\item $\begin{pmatrix} 0.88 & 1.52 \\ 1.52 & 0.88 \end{pmatrix}$
	\end{enumerate}
	
	\vspace{0.1cm}
	\textcolor{myblue}{\textbf{ANSWER: B}}
	
	\vspace{0.1cm}
	\textbf{DETAILED EXPLANATION:}
	\begin{itemize}
		\item Original covariance: $\begin{pmatrix} 1 & -0.58 \\ -0.58 & 1 \end{pmatrix}$
		\item Inverse: $\begin{pmatrix} 1.52 & 0.88 \\ 0.88 & 1.52 \end{pmatrix}$
		\item Used in discriminant function
		\item Features are standardized (variance = 1)
	\end{itemize}
	
	\vspace{0.1cm}
	\textcolor{myred}{\textbf{WHY OTHER OPTIONS ARE WRONG:}}
	\begin{itemize}
		\item \textcolor{myred}{A:} This is the covariance matrix
		\item \textcolor{myred}{C:} Wrong signs
		\item \textcolor{myred}{D:} Wrong arrangement
	\end{itemize}
	
	\textcolor{myred}{\textbf{EXAM TRAP:}} Inverse covariance matrix = \textcolor{myblue}{$\begin{pmatrix} 1.52 & 0.88 \\ 0.88 & 1.52 \end{pmatrix}$}!
\end{frame}
	
	\begin{frame}{Q84: LDA New Observation Classification - Tutorial}
		\fontsize{6.5}{7.5}\selectfont
		\textbf{QUESTION: In the LDA example, why is the new borrower classified as no default?}
		
		\vspace{0.1cm}
		\textbf{OPTIONS:}
		\begin{enumerate}
			\item Because the default score is higher
			\item Because the no-default score is higher (0.30 > -3.95)
			\item Because of random chance
			\item Because there are more non-defaults in training
		\end{enumerate}
		
		\vspace{0.1cm}
		\textcolor{myblue}{\textbf{ANSWER: B}}
		
		\vspace{0.1cm}
		\textbf{DETAILED EXPLANATION:}
		\begin{itemize}
			\item $\delta_0 = 0.30$ (no default)
			\item $\delta_1 = -3.95$ (default)
			\item $0.30 > -3.95$, so classify as no default
			\item Discriminant score much higher for no default
		\end{itemize}
		
		\vspace{0.1cm}
		\textcolor{myred}{\textbf{WHY OTHER OPTIONS ARE WRONG:}}
		\begin{itemize}
			\item \textcolor{myred}{A:} Default score is lower
			\item \textcolor{myred}{C:} Classification is deterministic
			\item \textcolor{myred}{D:} Priors help but scores determine
		\end{itemize}
		
		\textcolor{myred}{\textbf{EXAM TRAP:}} Higher discriminant score = \textcolor{myblue}{preferred classification}!
	\end{frame}
	
	\begin{frame}{Q85: LDA vs Logistic Regression Summary - Tutorial}
		\fontsize{6.5}{7.5}\selectfont
		\textbf{QUESTION: When is LDA preferred over logistic regression?}
		
		\vspace{0.1cm}
		\textbf{OPTIONS:}
		\begin{enumerate}
			\item For binary classification only
			\item For multi-class classification with well-separated classes
			\item For continuous prediction
			\item For clustering
		\end{enumerate}
		
		\vspace{0.1cm}
		\textcolor{myblue}{\textbf{ANSWER: B}}
		
		\vspace{0.1cm}
		\textbf{DETAILED EXPLANATION:}
		\begin{itemize}
			\item LDA handles multiple classes naturally
			\item Works well when classes are well-separated
			\item Logistic regression needs extension for multi-class
			\item LDA is robust even when assumptions violated
		\end{itemize}
		
		\vspace{0.1cm}
		\textcolor{myred}{\textbf{WHY OTHER OPTIONS ARE WRONG:}}
		\begin{itemize}
			\item \textcolor{myred}{A:} Logistic regression is for binary
			\item \textcolor{myred}{C:} Neither is for continuous
			\item \textcolor{myred}{D:} Neither is for clustering
		\end{itemize}
		
		\textcolor{myred}{\textbf{EXAM TRAP:}} LDA = \textcolor{myblue}{multi-class with well-separated classes}!
	\end{frame}
	
	% ============================================================
	% SECTION 7: CHAPTER QUESTIONS - QUESTIONS 86-100
	% ============================================================
	
	\begin{frame}{Q86: Outlier Detection Methods - Tutorial}
		\fontsize{6.5}{7.5}\selectfont
		\textbf{QUESTION: What are two methods to detect outliers in regression?}
		
		\vspace{0.1cm}
		\textbf{OPTIONS:}
		\begin{enumerate}
			\item Examine residuals and Cook's distance
			\item Examine R-squared and p-values
			\item Examine correlation and covariance
			\item Examine intercept and slope
		\end{enumerate}
		
		\vspace{0.1cm}
		\textcolor{myblue}{\textbf{ANSWER: A}}
		
		\vspace{0.1cm}
		\textbf{DETAILED EXPLANATION:}
		\begin{itemize}
			\item Residual plots: look for large residuals
			\item Cook's distance: measures influence of each point
			\item Large Cook's distance warrants investigation
			\item Both help identify influential points
		\end{itemize}
		
		\vspace{0.1cm}
		\textcolor{myred}{\textbf{WHY OTHER OPTIONS ARE WRONG:}}
		\begin{itemize}
			\item \textcolor{myred}{B:} R-squared and p-values don't detect outliers
			\item \textcolor{myred}{C:} Correlation doesn't detect outliers
			\item \textcolor{myred}{D:} Intercept/slope don't detect outliers
		\end{itemize}
		
		\textcolor{myred}{\textbf{EXAM TRAP:}} Outliers detected by \textcolor{myblue}{residuals and Cook's distance}!
	\end{frame}
	
	\begin{frame}{Q87: Stepwise Regression Explanation - Tutorial}
		\fontsize{6.5}{7.5}\selectfont
		\textbf{QUESTION: How does stepwise regression work?}
		
		\vspace{0.1cm}
		\textbf{OPTIONS:}
		\begin{enumerate}
			\item It randomly selects features
			\item It adds or removes features based on a criterion
			\item It always includes all features
			\item It never removes features
		\end{enumerate}
		
		\vspace{0.1cm}
		\textcolor{myblue}{\textbf{ANSWER: B}}
		
		\vspace{0.1cm}
		\textbf{DETAILED EXPLANATION:}
		\begin{itemize}
			\item Forward: start with no features, add best one each step
			\item Backward: start with all features, remove worst one each step
			\item Uses AIC or similar criterion
			\item Stops when no improvement
		\end{itemize}
		
		\vspace{0.1cm}
		\textcolor{myred}{\textbf{WHY OTHER OPTIONS ARE WRONG:}}
		\begin{itemize}
			\item \textcolor{myred}{A:} Stepwise is systematic, not random
			\item \textcolor{myred}{C:} Stepwise selects subset, not all
			\item \textcolor{myred}{D:} Stepwise can remove features
		\end{itemize}
		
		\textcolor{myred}{\textbf{EXAM TRAP:}} Stepwise = \textcolor{myblue}{systematic add/remove based on criterion}!
	\end{frame}
	
	\begin{frame}{Q88: Linear Regression for Binary Data - Tutorial}
		\fontsize{6.5}{7.5}\selectfont
		\textbf{QUESTION: Why can't linear regression be used for binary outcomes?}
		
		\vspace{0.1cm}
		\textbf{OPTIONS:}
		\begin{enumerate}
			\item Because it's too complex
			\item Because predictions can fall outside 0-1 range
			\item Because it always gives 0.5
			\item Because it requires categorical features
		\end{enumerate}
		
		\vspace{0.1cm}
		\textcolor{myblue}{\textbf{ANSWER: B}}
		
		\vspace{0.1cm}
		\textbf{DETAILED EXPLANATION:}
		\begin{itemize}
			\item Linear regression gives unbounded predictions
			\item Could predict negative probabilities
			\item Could predict probabilities > 1
			\item Logistic regression bounds output between 0 and 1
		\end{itemize}
		
		\vspace{0.1cm}
		\textcolor{myred}{\textbf{WHY OTHER OPTIONS ARE WRONG:}}
		\begin{itemize}
			\item \textcolor{myred}{A:} Linear regression is simpler
			\item \textcolor{myred}{C:} Linear regression doesn't always give 0.5
			\item \textcolor{myred}{D:} Linear regression doesn't require categorical features
		\end{itemize}
		
		\textcolor{myred}{\textbf{EXAM TRAP:}} Linear regression = \textcolor{myblue}{unbounded predictions} (can be outside 0-1)!
	\end{frame}
	
	\begin{frame}{Q89: LDA Assumptions Detail - Tutorial}
		\fontsize{6.5}{7.5}\selectfont
		\textbf{QUESTION: What are the main assumptions underlying LDA?}
		
		\vspace{0.1cm}
		\textbf{OPTIONS:}
		\begin{enumerate}
			\item Features are normally distributed with common covariance
			\item Features are independent and identically distributed
			\item Features are uniformly distributed
			\item Features are categorical
		\end{enumerate}
		
		\vspace{0.1cm}
		\textcolor{myblue}{\textbf{ANSWER: A}}
		
		\vspace{0.1cm}
		\textbf{DETAILED EXPLANATION:}
		\begin{itemize}
			\item Multivariate normal for each class
			\item Different mean vectors per class
			\item Common covariance matrix
			\item Features must be continuous
		\end{itemize}
		
		\vspace{0.1cm}
		\textcolor{myred}{\textbf{WHY OTHER OPTIONS ARE WRONG:}}
		\begin{itemize}
			\item \textcolor{myred}{B:} Features need not be independent
			\item \textcolor{myred}{C:} Normal, not uniform
			\item \textcolor{myred}{D:} Features are continuous, not categorical
		\end{itemize}
		
		\textcolor{myred}{\textbf{EXAM TRAP:}} LDA = \textcolor{myblue}{normal distributions + common covariance}!
	\end{frame}
	
	\begin{frame}{Q90: ROA and SIZE Example - Tutorial}
		\fontsize{6.5}{7.5}\selectfont
		\textbf{QUESTION: In the ROA and SIZE example, is the relationship linear?}
		
		\vspace{0.1cm}
		\textbf{OPTIONS:}
		\begin{enumerate}
			\item Yes, it's perfectly linear
			\item No, because the equation includes a squared term
			\item Yes, because all coefficients are positive
			\item No, because ROA is in percent
		\end{enumerate}
		
		\vspace{0.1cm}
		\textcolor{myblue}{\textbf{ANSWER: B}}
		
		\vspace{0.1cm}
		\textbf{DETAILED EXPLANATION:}
		\begin{itemize}
			\item Equation: $ROA = 2.6 + 4.22(SIZE) - 1.32(SIZE)^2$
			\item Squared term creates non-linearity
			\item Relationship changes with SIZE
			\item Maximum at SIZE = 1.60
		\end{itemize}
		
		\vspace{0.1cm}
		\textcolor{myred}{\textbf{WHY OTHER OPTIONS ARE WRONG:}}
		\begin{itemize}
			\item \textcolor{myred}{A:} Squared term makes it non-linear
			\item \textcolor{myred}{C:} Positive coefficients don't imply linearity
			\item \textcolor{myred}{D:} Units don't determine linearity
		\end{itemize}
		
		\textcolor{myred}{\textbf{EXAM TRAP:}} Squared term = \textcolor{myblue}{non-linear relationship}!
	\end{frame}
	
	\begin{frame}{Q91: ROA Calculation Example - Tutorial}
		\fontsize{6.5}{7.5}\selectfont
		\textbf{QUESTION: For the ROA model, what is the expected ROA for SIZE = 1.0?}
		
		\vspace{0.1cm}
		\textbf{OPTIONS:}
		\begin{enumerate}
			\item 4.50\%
			\item 5.50\%
			\item 6.50\%
			\item 7.50\%
		\end{enumerate}
		
		\vspace{0.1cm}
		\textcolor{myblue}{\textbf{ANSWER: B}}
		
		\vspace{0.1cm}
		\textbf{DETAILED EXPLANATION:}
		\begin{itemize}
			\item Equation: $ROA = 2.6 + 4.22(1.0) - 1.32(1.0)^2$
			\item $ROA = 2.6 + 4.22 - 1.32$
			\item $ROA = 5.50$
			\item SIZE measured in hundreds of millions
		\end{itemize}
		
		\vspace{0.1cm}
		\textcolor{myred}{\textbf{WHY OTHER OPTIONS ARE WRONG:}}
		\begin{itemize}
			\item \textcolor{myred}{A:} Too low
			\item \textcolor{myred}{C:} Too high
			\item \textcolor{myred}{D:} Too high
		\end{itemize}
		
		\textcolor{myred}{\textbf{EXAM TRAP:}} Plug SIZE = 1.0 into equation: \textcolor{myblue}{ROA = 5.50\%}!
	\end{frame}
	
	\begin{frame}{Q92: Optimal SIZE Calculation - Tutorial}
		\fontsize{6.5}{7.5}\selectfont
		\textbf{QUESTION: In the ROA example, what SIZE maximizes ROA?}
		
		\vspace{0.1cm}
		\textbf{OPTIONS:}
		\begin{enumerate}
			\item 1.00
			\item 1.60
			\item 2.00
			\item 0.50
		\end{enumerate}
		
		\vspace{0.1cm}
		\textcolor{myblue}{\textbf{ANSWER: B}}
		
		\vspace{0.1cm}
		\textbf{DETAILED EXPLANATION:}
		\begin{itemize}
			\item Derivative: $\frac{dROA}{dSIZE} = 4.22 - 2.64(SIZE)$
			\item Set to zero: $4.22 - 2.64(SIZE) = 0$
			\item $SIZE = 4.22 / 2.64 = 1.60$
			\item This is \$160 million (since SIZE is in hundreds of millions)
		\end{itemize}
		
		\vspace{0.1cm}
		\textcolor{myred}{\textbf{WHY OTHER OPTIONS ARE WRONG:}}
		\begin{itemize}
			\item \textcolor{myred}{A:} Not the maximum
			\item \textcolor{myred}{C:} Beyond the maximum
			\item \textcolor{myred}{D:} Below the maximum
		\end{itemize}
		
		\textcolor{myred}{\textbf{EXAM TRAP:}} Optimal SIZE = \textcolor{myblue}{1.60} (hundreds of millions)!
	\end{frame}
	
	\begin{frame}{Q93: Maximum ROA Calculation - Tutorial}
		\fontsize{6.5}{7.5}\selectfont
		\textbf{QUESTION: What is the maximum ROA in the example?}
		
		\vspace{0.1cm}
		\textbf{OPTIONS:}
		\begin{enumerate}
			\item 5.50\%
			\item 6.03\%
			\item 7.00\%
			\item 8.50\%
		\end{enumerate}
		
		\vspace{0.1cm}
		\textcolor{myblue}{\textbf{ANSWER: B}}
		
		\vspace{0.1cm}
		\textbf{DETAILED EXPLANATION:}
		\begin{itemize}
			\item Plug SIZE = 1.60 into equation
			\item $ROA = 2.6 + 4.22(1.60) - 1.32(1.60)^2$
			\item $ROA = 2.6 + 6.752 - 3.3792$
			\item $ROA = 6.03\%$
		\end{itemize}
		
		\vspace{0.1cm}
		\textcolor{myred}{\textbf{WHY OTHER OPTIONS ARE WRONG:}}
		\begin{itemize}
			\item \textcolor{myred}{A:} This is ROA at SIZE = 1.0
			\item \textcolor{myred}{C:} Too high
			\item \textcolor{myred}{D:} Too high
		\end{itemize}
		
		\textcolor{myred}{\textbf{EXAM TRAP:}} Maximum ROA = \textcolor{myblue}{6.03\%} at SIZE = 1.60!
	\end{frame}
	
	\begin{frame}{Q94: Marketing Campaign LDA - Tutorial}
		\fontsize{6.5}{7.5}\selectfont
		\textbf{QUESTION: In the marketing campaign example, what are the features?}
		
		\vspace{0.1cm}
		\textbf{OPTIONS:}
		\begin{enumerate}
			\item Age and Income
			\item Balance and Age
			\item Balance and Income
			\item Age and Gender
		\end{enumerate}
		
		\vspace{0.1cm}
		\textcolor{myblue}{\textbf{ANSWER: B}}
		
		\vspace{0.1cm}
		\textbf{DETAILED EXPLANATION:}
		\begin{itemize}
			\item Features: Balance (dollars) and Age (years)
			\item Target: Subscribed (Yes/No)
			\item Binary classification problem
			\item LDA used to classify clients
		\end{itemize}
		
		\vspace{0.1cm}
		\textcolor{myred}{\textbf{WHY OTHER OPTIONS ARE WRONG:}}
		\begin{itemize}
			\item \textcolor{myred}{A:} Income is not a feature
			\item \textcolor{myred}{C:} Income is not a feature
			\item \textcolor{myred}{D:} Gender is not a feature
		\end{itemize}
		
		\textcolor{myred}{\textbf{EXAM TRAP:}} Features = \textcolor{myblue}{Balance and Age}!
	\end{frame}
	
	\begin{frame}{Q95: LDA Prior in Marketing Example - Tutorial}
		\fontsize{6.5}{7.5}\selectfont
		\textbf{QUESTION: In the marketing example, what are the prior probabilities?}
		
		\vspace{0.1cm}
		\textbf{OPTIONS:}
		\begin{enumerate}
			\item 0.5 and 0.5
			\item 0.6 and 0.4
			\item 0.7 and 0.3
			\item 0.8 and 0.2
		\end{enumerate}
		
		\vspace{0.1cm}
		\textcolor{myblue}{\textbf{ANSWER: B}}
		
		\vspace{0.1cm}
		\textbf{DETAILED EXPLANATION:}
		\begin{itemize}
			\item 6 out of 10 subscribed (60\%)
			\item 4 out of 10 did not subscribe (40\%)
			\item $\pi_1 = 0.6$ (subscribed)
			\item $\pi_2 = 0.4$ (not subscribed)
		\end{itemize}
		
		\vspace{0.1cm}
		\textcolor{myred}{\textbf{WHY OTHER OPTIONS ARE WRONG:}}
		\begin{itemize}
			\item \textcolor{myred}{A:} Not balanced (6 vs 4)
			\item \textcolor{myred}{C:} Not 7 vs 3
			\item \textcolor{myred}{D:} Not 8 vs 2
		\end{itemize}
		
		\textcolor{myred}{\textbf{EXAM TRAP:}} Priors = \textcolor{myblue}{0.6 and 0.4} from data!
	\end{frame}
	
	\begin{frame}{Q96: Marketing Example Class Means - Tutorial}
		\fontsize{6.5}{7.5}\selectfont
		\textbf{QUESTION: In the marketing example, what is the mean balance for subscribers?}
		
		\vspace{0.1cm}
		\textbf{OPTIONS:}
		\begin{enumerate}
			\item -0.44
			\item 0.29
			\item 0.00
			\item 1.00
		\end{enumerate}
		
		\vspace{0.1cm}
		\textcolor{myblue}{\textbf{ANSWER: B}}
		
		\vspace{0.1cm}
		\textbf{DETAILED EXPLANATION:}
		\begin{itemize}
			\item $\mu_1 = \begin{pmatrix} 0.29 \\ 0.00 \end{pmatrix}$
			\item First element: mean balance for subscribers
			\item Second element: mean age for subscribers
			\item Features are standardized
		\end{itemize}
		
		\vspace{0.1cm}
		\textcolor{myred}{\textbf{WHY OTHER OPTIONS ARE WRONG:}}
		\begin{itemize}
			\item \textcolor{myred}{A:} This is mean balance for non-subscribers
			\item \textcolor{myred}{C:} This is mean age for subscribers
			\item \textcolor{myred}{D:} Not correct
		\end{itemize}
		
		\textcolor{myred}{\textbf{EXAM TRAP:}} Subscriber balance mean = \textcolor{myblue}{0.29} (standardized)!
	\end{frame}
	
	\begin{frame}{Q97: Marketing Example Non-Subscriber Means - Tutorial}
		\fontsize{6.5}{7.5}\selectfont
		\textbf{QUESTION: In the marketing example, what is the mean balance for non-subscribers?}
		
		\vspace{0.1cm}
		\textbf{OPTIONS:}
		\begin{enumerate}
			\item 0.29
			\item -0.44
			\item 0.00
			\item -1.00
		\end{enumerate}
		
		\vspace{0.1cm}
		\textcolor{myblue}{\textbf{ANSWER: B}}
		
		\vspace{0.1cm}
		\textbf{DETAILED EXPLANATION:}
		\begin{itemize}
			\item $\mu_2 = \begin{pmatrix} -0.44 \\ 0.00 \end{pmatrix}$
			\item First element: mean balance for non-subscribers
			\item Second element: mean age for non-subscribers
			\item Non-subscribers have lower average balance
		\end{itemize}
		
		\vspace{0.1cm}
		\textcolor{myred}{\textbf{WHY OTHER OPTIONS ARE WRONG:}}
		\begin{itemize}
			\item \textcolor{myred}{A:} This is mean balance for subscribers
			\item \textcolor{myred}{C:} This is mean age
			\item \textcolor{myred}{D:} Not correct
		\end{itemize}
		
		\textcolor{myred}{\textbf{EXAM TRAP:}} Non-subscriber balance mean = \textcolor{myblue}{-0.44} (standardized)!
	\end{frame}
	
	\begin{frame}{Q98: Marketing Example Covariance - Tutorial}
		\fontsize{6.5}{7.5}\selectfont
		\textbf{QUESTION: In the marketing example, what is the covariance between features?}
		
		\vspace{0.1cm}
		\textbf{OPTIONS:}
		\begin{enumerate}
			\item 0.33
			\item 0.58
			\item -0.58
			\item 1.00
		\end{enumerate}
		
		\vspace{0.1cm}
		\textcolor{myblue}{\textbf{ANSWER: A}}
		
		\vspace{0.1cm}
		\textbf{DETAILED EXPLANATION:}
		\begin{itemize}
			\item Covariance matrix: $\begin{pmatrix} 1 & 0.33 \\ 0.33 & 1 \end{pmatrix}$
			\item Variances = 1 (features standardized)
			\item Covariance = 0.33
			\item Positive correlation between balance and age
		\end{itemize}
		
		\vspace{0.1cm}
		\textcolor{myred}{\textbf{WHY OTHER OPTIONS ARE WRONG:}}
		\begin{itemize}
			\item \textcolor{myred}{B:} This is from LDA example
			\item \textcolor{myred}{C:} This is from LDA example
			\item \textcolor{myred}{D:} This is variance, not covariance
		\end{itemize}
		
		\textcolor{myred}{\textbf{EXAM TRAP:}} Covariance = \textcolor{myblue}{0.33} in marketing example!
	\end{frame}
	
	\begin{frame}{Q99: Marketing Example Classification - Tutorial}
		\fontsize{6.5}{7.5}\selectfont
		\textbf{QUESTION: In the marketing example, how would a 51-year-old with \$10,635 be classified?}
		
		\vspace{0.1cm}
		\textbf{OPTIONS:}
		\begin{enumerate}
			\item Subscribed
			\item Not subscribed
			\item Cannot be determined
			\item Both equally likely
		\end{enumerate}
		
		\vspace{0.1cm}
		\textcolor{myblue}{\textbf{ANSWER: A}}
		
		\vspace{0.1cm}
		\textbf{DETAILED EXPLANATION:}
		\begin{itemize}
			\item Standardized values: Age = 0.65, Balance = 1.64
			\item LDA would classify based on discriminant scores
			\item Higher balance and older age suggest subscription
			\item Subscriber class has higher mean balance
		\end{itemize}
		
		\vspace{0.1cm}
		\textcolor{myred}{\textbf{WHY OTHER OPTIONS ARE WRONG:}}
		\begin{itemize}
			\item \textcolor{myred}{B:} Features suggest subscriber
			\item \textcolor{myred}{C:} We have enough information
			\item \textcolor{myred}{D:} Scores would not be equal
		\end{itemize}
		
		\textcolor{myred}{\textbf{EXAM TRAP:}} Classification uses \textcolor{myblue}{discriminant scores} to decide!
	\end{frame}
	
	\begin{frame}{Q100: Final Summary - Tutorial}
		\fontsize{6.5}{7.5}\selectfont
		\textbf{QUESTION: What is the most comprehensive summary of supervised learning for numerical data?}
		
		\vspace{0.1cm}
		\textbf{OPTIONS:}
		\begin{enumerate}
			\item Only linear regression is used
			\item Models include linear regression (continuous outcomes) and logistic regression/LDA (classification)
			\item Only logistic regression is used
			\item Only LDA is used
		\end{enumerate}
		
		\vspace{0.1cm}
		\textcolor{myblue}{\textbf{ANSWER: B}}
		
		\vspace{0.1cm}
		\textbf{DETAILED EXPLANATION:}
		\begin{itemize}
			\item \textbf{Continuous outcomes:} Linear regression (simple, multiple)
			\item \textbf{Classification:} Logistic regression, LDA
			\item \textbf{Extensions:} Multinomial logit, ordered logit
			\item \textbf{Key concepts:} OLS, AIC, stepwise, assumptions
		\end{itemize}
		
		\vspace{0.1cm}
		\textcolor{myred}{\textbf{WHY OTHER OPTIONS ARE WRONG:}}
		\begin{itemize}
			\item \textcolor{myred}{A:} Multiple techniques are used
			\item \textcolor{myred}{C:} Linear regression and LDA also used
			\item \textcolor{myred}{D:} Linear and logistic also used
		\end{itemize}
		
		\textcolor{myred}{\textbf{EXAM SUCCESS:}} Understand \textcolor{myblue}{regression and classification} techniques!
	\end{frame}
	
\end{document}